\documentclass[11pt]{article}

\usepackage{amsmath,amssymb,mathtools,bm,booktabs,longtable}
\usepackage[margin=1in]{geometry}

\newcommand{\dd}{\,\mathrm d}
\newcommand{\R}{\mathbb R}
\newcommand{\E}{\mathbb E}
\newcommand{\Tr}{\operatorname{Tr}}
\newcommand{\diag}{\operatorname{diag}}
\newcommand{\rank}{\operatorname{rank}}
\newcommand{\argmax}{\operatorname*{arg\,max}}
\newcommand{\argmin}{\operatorname*{arg\,min}}
\newcommand{\one}{\mathbf 1}
\newcommand{\zero}{\mathbf 0}
\newcommand{\eps}{\varepsilon}

\begin{document}

\section*{Complete Projected Analytical Conductivity Model}
\subsection*{From Composition to Generator to Scalar Conductivity}

\subsection*{Purpose}

The model maps one electrolyte composition \(x\) to one scalar conductivity
\[
\sigma(x)\quad [\mathrm{S\,m^{-1}}],
\]
or equivalently
\[
\sigma_{\mathrm{mS/cm}}(x)=10\,\sigma(x).
\]

No Green--Kubo time integral is evaluated in production.
No Einstein--Helfand long-time slope is evaluated in production.
The model computes conductivity from projected equilibrium, flux, displacement,
self-current, and memory primitives.

\[
x
\longmapsto
\mathcal L
\longmapsto
(U_x,D_x,P_x,\mathcal A,\Psi)
\longmapsto
(c,K,Q,d,M,D^{\rm self},A,h)
\longmapsto
\sigma.
\]

Here \(\mathcal L\) is the physical species library (defined below) that turns
a bare recipe into the numerical functions \(U_x,D_x,P_x,\mathcal A,\Psi\)
required by the generator formalism. The document is organized so that every
symbol appearing in the final boxed formula is traceable, in order, back to
the composition \(x\).

\section{Notation and operation glossary}

\subsection{Vectors, matrices, tensors}

A vector is a column vector unless stated otherwise.

For vectors \(u,v\in\R^n\),
\[
u^\top v=\sum_{i=1}^n u_i v_i.
\]

For a square matrix \(B\),
\[
\Tr(B)=\sum_i B_{ii}.
\]

For a symmetric matrix \(B\),
\[
B\succeq0
\]
means
\[
y^\top B y\ge0
\qquad
\text{for all }y.
\]

For a tensor \(T_i\in\R^{3\times3}\),
\[
\Tr(T_i)=T_{i,xx}+T_{i,yy}+T_{i,zz}.
\]

\subsection{The symbol \(B^\#\)}

The symbol
\[
B^\#
\]
means the Moore--Penrose pseudoinverse when \(B\) is a symmetric positive
semidefinite matrix.

Implementation:

1. Compute the eigendecomposition
\[
B=V\Lambda V^\top,
\]
where
\[
\Lambda=\diag(\lambda_1,\ldots,\lambda_m),
\qquad
\lambda_r\ge0.
\]

2. Choose numerical tolerance
\[
\tau_{\#}=10^{-12}\max_r \lambda_r.
\]

3. Define
\[
\lambda_r^\#=
\begin{cases}
1/\lambda_r, & \lambda_r>\tau_{\#},\\
0, & \lambda_r\le\tau_{\#}.
\end{cases}
\]

4. Return
\[
B^\#=V\diag(\lambda_1^\#,\ldots,\lambda_m^\#)V^\top.
\]

If a vector \(y\) has nonzero projection into the nullspace of \(B\), then
\[
y^\top B^\# y
\]
silently discards that null component. For a physical memory matrix this is
allowed only if the discarded component is numerically zero. Therefore require
\[
\|(I-BB^\#)y\| \le \tau_{\rm null}\|y\|
\]
for each current-coupling vector \(y\), with for example
\[
\tau_{\rm null}=10^{-8}.
\]

For the finite Markov generator \(Q\), do not directly form \((-Q)^\#\).
Instead solve the constrained Poisson system defined in Section~\ref{sec:finitestate}.

\section{Composition input}
\label{sec:compinput}

The electrolyte composition is
\[
x=
\left(
T,\,
\{C_a\}_{a\in\mathcal C},\,
\{\phi_s\}_{s\in\mathcal S},\,
\mathcal P
\right).
\]

Definitions:

\[
T>0
\]
is temperature in K.

\[
\mathcal C
\]
is the set of conserved components, for example Li, PF6, FSI, solvent molecules,
neutral additives, or coarse-grained site classes.

\[
C_a
\quad [\mathrm{mol\,m^{-3}}]
\]
is the analytical total concentration of conserved component \(a\).

\[
\phi_s
\]
is the liquid fraction of solvent/additive species \(s\).

\[
\mathcal P
\]
is the physical parameter map needed to construct the generator:
\[
\mathcal P
=
\left(
\text{charge sites},
\text{topology},
\text{interaction parameters},
\text{mobility parameters},
\text{basin definitions},
\text{memory coordinate definitions}
\right).
\]

A recipe containing only names such as EC, DMC, LiPF6, FEC, VC, FSI, and
concentrations is not mathematically enough. The model requires enough
information to build the following functions:
\[
U_x(q),
\qquad
D_x(q),
\qquad
P_x(q),
\qquad
A_i\subset\Omega_x,
\qquad
\psi_\alpha(q).
\]

\(\mathcal P\) as stated above is an input \emph{requirement}, not yet a
\emph{construction}. Section~\ref{sec:c2g} makes this precise: it shows why
\(\mathcal P\) cannot be inferred from species names and concentrations alone,
and Section~\ref{sec:library} expands \(\mathcal P\) into an explicit,
buildable species library \(\mathcal L\) from which every one of
\(U_x,D_x,P_x,A_i,\psi_\alpha\) is derived mechanically.

\section{Composition unit conversion}
\label{sec:unitconv}

The model accepts electrolyte formulations as they appear in experimental
practice:
\begin{itemize}
\item Solvents in volume/volume fractions (v/v),
\item Salts in molar concentration (mol~L$^{-1}$, i.e.\ M),
\item Additives in weight percent (wt\%).
\end{itemize}

This section converts those external units into the canonical model inputs
\(\{C_a\}_{a\in\mathcal C}\) in \(\mathrm{mol\,m^{-3}}\) and
\(\{\phi_s\}_{s\in\mathcal S}\) in dimensionless volume fractions.

\subsection{Temperature}

Temperature is supplied in \(\mathrm K\). If supplied in \({}^\circ\mathrm C\):
\[
T = T_{{}^\circ\mathrm C} + 273.15.
\]

\subsection{Solvents: volume fraction}

Suppose the formulation specifies solvents \(s_1,\ldots,s_{n_{\rm solv}}\) at
volume-to-volume ratios \(v_1:v_2:\cdots:v_{n_{\rm solv}}\).

Step 1. Normalize to volume fractions:
\[
\phi_{s_k}^{\rm raw}
=
\frac{v_k}{\sum_{\ell=1}^{n_{\rm solv}} v_\ell}.
\]

Step 2. Correct for additive volume displacement. Let the additive mass
fractions \(f_{\rm add,j}^{\rm wt}\) (defined below) add a total volume
\(V_{\rm add}\) per unit of total system volume. Set
\[
\phi_{s_k}
=
\phi_{s_k}^{\rm raw}\,(1-V_{\rm add}).
\]

For the initial computation, set \(V_{\rm add}=0\) and update after additives
are converted, iterating once or until convergence.

Step 3. Pass \(\phi_{s_k}\) to the molecule-count algorithm
(Section~\ref{sec:library}).

\subsection{Salts: molarity to SI concentration}

A salt concentration specified in \(\mathrm{mol\,L^{-1}}\) (molarity) is
converted to \(\mathrm{mol\,m^{-3}}\) (SI) by:
\[
C_{\rm salt}^{\rm SI}
=
C_{\rm salt}^{\rm M} \times 10^3
\quad [\mathrm{mol\,m^{-3}}].
\]

For a salt \(\mathrm{Li_{\nu_+}A_{\nu_-}}\) with stoichiometric coefficients
\(\nu_+,\nu_-\):
\[
C_{\rm Li} = \nu_+ C_{\rm salt}^{\rm SI},
\qquad
C_{\rm A} = \nu_- C_{\rm salt}^{\rm SI}.
\]

For mixed salts (e.g.\ \(0.6\,\mathrm{M}\) LiPF$_6$ \(+\) \(0.4\,\mathrm{M}\)
LiFSI), all sources contribute to the same Li pool:
\[
C_{\rm Li}
=
\sum_{\rm salts} \nu_{+,k} C_{\rm salt\,k}^{\rm SI}.
\]

Separate anion concentrations are kept distinct:
\[
C_{\rm PF6} = \nu_{-,\rm PF6} C_{\rm LiPF6}^{\rm SI},
\qquad
C_{\rm FSI} = \nu_{-,\rm FSI} C_{\rm LiFSI}^{\rm SI}.
\]

Charge neutrality check:
\[
\sum_a z_a^{\rm net} C_a = 0.
\]

If this fails, the formulation is physically inconsistent.

\subsection{Additives: weight percent to concentration}

An additive \(j\) specified in wt\% is converted as follows.

Let:
\[
f_j^{\rm wt} = \frac{\text{mass of additive }j}{\text{total solution mass}}
\]
be the weight fraction (wt\%$/100$).

The solution density at composition \(x\) is estimated as a mixture rule:
\[
\rho_{\rm soln}
=
\frac{1}
{\displaystyle
\sum_s \phi_s/\rho_s
+
\sum_{\rm salts} C_k^{\rm SI} M_k / (\rho_k N_A \times 10^3)
},
\]
where \(\rho_s\) is pure solvent density in \(\mathrm{kg\,m^{-3}}\), \(M_k\)
is salt molar mass in \(\mathrm{g\,mol^{-1}}\), and the sum for salts accounts
for the dissolved ion volume. An iterative estimate starting from solvent-only
density converges in 1--3 steps.

Then mass concentration is:
\[
\rho_j^{\rm mass}
=
f_j^{\rm wt}\,\rho_{\rm soln}
\quad [\mathrm{kg\,m^{-3}}].
\]

Molar concentration:
\[
C_j
=
\frac{\rho_j^{\rm mass}}{M_j\times 10^{-3}}
\quad [\mathrm{mol\,m^{-3}}],
\]
where \(M_j\) is the additive molar mass in \(\mathrm{g\,mol^{-1}}\).

If the additive is a \emph{neutral ligand} (e.g.\ FEC, VC) and its
concentration is small relative to solvents, it may optionally be expressed as
a volume fraction instead:
\[
\phi_j
=
\frac{C_j M_j \times 10^{-3}}{\rho_j}
\quad [\text{dimensionless}].
\]

The model accepts \emph{both}: include the additive in \(\{C_a\}\) with a
stoichiometry vector \(\nu_{i,j}\) if it participates in basin/motif formation,
or in \(\{\phi_s\}\) if it acts purely as a dilute solvent. Do not count it
twice.

\subsection{Canonical input tuple}

After the conversions above, the canonical model input is
\[
x
=
\bigl(
T,\,
\{C_a\}_{a\in\mathcal C},\,
\{\phi_s\}_{s\in\mathcal S},\,
\mathcal L
\bigr),
\]
with \(C_a\) in \(\mathrm{mol\,m^{-3}}\), \(\phi_s\) dimensionless, and
\(\mathcal L\) the physical library (Section~\ref{sec:library}).

\subsection{Summary conversion table}

\[
\begin{array}{lll}
\toprule
\text{External unit} & \text{Symbol} & \text{Canonical model input} \\
\midrule
\text{Solvent v/v ratio } v_k & v_1:\cdots:v_{n} &
\phi_{s_k}=v_k/\sum_\ell v_\ell \\[4pt]
\text{Salt molarity} & C^{\rm M}\,[\mathrm{M}] &
C_a^{\rm SI}=C^{\rm M}\times 10^3\,[\mathrm{mol\,m^{-3}}] \\[4pt]
\text{Additive wt\%} & f^{\rm wt}\,[\%] &
C_j = f^{\rm wt}\rho_{\rm soln}/(100\,M_j\times 10^{-3})\,[\mathrm{mol\,m^{-3}}] \\[4pt]
\text{Temperature} & T_{{}^\circ\mathrm C}\,[{}^\circ\mathrm C] &
T=T_{{}^\circ\mathrm C}+273.15\,[\mathrm K] \\
\bottomrule
\end{array}
\]

\section{Composition-to-generator construction}
\label{sec:c2g}

\subsection{Why recipe names alone are not a mathematical model}

Let a recipe-only composition be
\[
r=
\left(
T,\,
\{C_a\}_{a\in\mathcal C},\,
\{\phi_s\}_{s\in\mathcal S}
\right),
\]
where \(C_a\) are analytical concentrations and \(\phi_s\) are liquid fractions.
This is exactly the tuple \((T,\{C_a\},\{\phi_s\})\) of Section~\ref{sec:compinput},
with \(\mathcal P\) omitted.

A recipe-only object \(r\) does not determine conductivity, because it does
not determine:
\[
U_x(q),\qquad D_x(q),\qquad P_x(q),\qquad A_i,\qquad \psi_\alpha(q).
\]

Therefore the conductivity map is not
\[
r\mapsto \sigma.
\]

The full composition is
\[
x=(r,\mathcal L),
\]
where \(\mathcal L\) is the physical species library. The conductivity map is
\[
(r,\mathcal L)\mapsto \sigma.
\]

If \(\mathcal L\) is absent, the first-principles map is undefined. Equivalently,
\(\mathcal P=\mathcal P(\mathcal L)\): the physical parameter map of
Section~\ref{sec:compinput} is precisely what the library \(\mathcal L\)
supplies.

\subsection{Non-identifiability lemma}

Suppose only \(r\) is given. Let two libraries \(\mathcal L_1,\mathcal L_2\)
have the same species names and concentrations but different microscopic
mobility tensors:
\[
D_x^{(2)}(q)=\alpha D_x^{(1)}(q),
\qquad
\alpha\neq1.
\]

Keeping the same \(U_x,P_x,A_i,\psi_\alpha\), the self-current tensors scale as
\[
D_{i,2}^{\rm self}
=
\alpha D_{i,1}^{\rm self},
\]
and the capacity fluxes scale as
\[
K_{ij}^{(2)}
=
\alpha K_{ij}^{(1)}
\]
for overdamped Smoluchowski dynamics because
\[
K_{ij}
=
\int
\nabla q_{ij}^{\top}D_x\nabla q_{ij}\,\rho_x(\dd q).
\]

Thus, unless the memory correction exactly compensates for all scaling,
\[
\sigma(r,\mathcal L_2)\neq\sigma(r,\mathcal L_1).
\]

So \(r\) alone cannot uniquely define \(\sigma\). The library \(\mathcal L\) is
part of the mathematical input, not an implementation detail.

\section{Physical species library}
\label{sec:library}

The species library is
\[
\mathcal L=
\left(
\{\mathcal T_s\}_{s\in\mathcal S_{\rm mol}},
\{\mathcal Q_s\}_{s\in\mathcal S_{\rm mol}},
\{\Theta_{st}\}_{s,t},
\mathcal B,
\mathcal M,
\mathcal R,
\mathcal \Psi
\right).
\]

Here:
\[
\mathcal S_{\rm mol}
=
\text{set of molecular, ionic, solvent, salt, and additive species}.
\]
\[
\mathcal T_s
=
\text{topology of species }s.
\]
\[
\mathcal Q_s
=
\text{charge-site model of species }s.
\]
\[
\Theta_{st}
=
\text{interaction parameters between species }s,t.
\]
\[
\mathcal B
=
\text{basin-definition library}.
\]
\[
\mathcal M
=
\text{mobility/friction library}.
\]
\[
\mathcal R
=
\text{reaction-coordinate and transition-surface library}.
\]
\[
\mathcal \Psi
=
\text{Mori memory-coordinate library}.
\]

\subsection{Species topology}

For each species \(s\), define a finite set of sites
\[
\mathcal I_s=\{1,\ldots,n_s\}.
\]

Each site \(u\in\mathcal I_s\) has:
\[
m_{su}>0
\quad
\text{mass},
\]
\[
a_{su}>0
\quad
\text{hydrodynamic radius},
\]
\[
\sigma_{su}>0,\quad \epsilon_{su}\ge0
\quad
\text{short-range nonbonded parameters},
\]
\[
z_{su}\in\mathbb R
\quad
\text{formal or partial charge number},
\]
\[
\alpha_{su}\ge0
\quad
\text{polarizability, if used}.
\]

The topology also contains bonded edges:
\[
\mathcal E_s^{\rm bond}
\subset
\mathcal I_s\times\mathcal I_s,
\]
angles:
\[
\mathcal E_s^{\rm angle}
\subset
\mathcal I_s^3,
\]
and torsions:
\[
\mathcal E_s^{\rm torsion}
\subset
\mathcal I_s^4.
\]

\subsection{Molecule counts}

For a simulation cell of volume \(V_{\rm cell}\), component count is
\[
N_a=\operatorname{round}(C_aN_A V_{\rm cell}),
\]
where \(N_A\) is Avogadro's number.

For solvents/additives specified by volume fraction \(\phi_s\), with molar mass
\(M_s\) and density \(\rho_s\), the number density is
\[
n_s^{\rm pure}
=
\frac{\rho_s}{M_s}N_A.
\]

The target count is
\[
N_s
=
\operatorname{round}\left(
\phi_s n_s^{\rm pure} V_{\rm cell}
\right).
\]

If both mass/volume fractions and molar concentrations are specified, solve the
volume-filling constraint:
\[
V_{\rm cell}
=
\sum_s \frac{N_sM_s}{\rho_sN_A}
+
\sum_a \bar v_a N_a,
\]
where \(\bar v_a\) is the partial molar volume of component \(a\). This is solved
by fixed-point iteration:
\[
V_{\rm cell}^{(k+1)}
=
\sum_s \frac{N_s^{(k)}M_s}{\rho_sN_A}
+
\sum_a \bar v_a N_a^{(k)}.
\]

Stop when
\[
\frac{|V_{\rm cell}^{(k+1)}-V_{\rm cell}^{(k)}|}
{V_{\rm cell}^{(k)}}<\tau_V.
\]

\section{Configuration space}
\label{sec:configspace}

\subsection{Abstract coordinate}

Let
\[
q\in\Omega_x\subseteq\R^d
\]
be the model coordinate.

For atomistic coordinates:
\[
q=(r_1,\ldots,r_N),
\qquad
r_a\in\R^3.
\]

For a reduced analytical model:
\[
q=
\left(
r_{\rm Li-A},
s_{\rm pair},
n_{\rm Li-solv},
n_{\rm Li-lig},
\omega_{\rm A},
n_{\rm cluster},
\rho_{\rm local},
p_{\rm atm},
\ldots
\right).
\]

Every component of \(q\) must have:

1. a unit,
2. a physical meaning,
3. a gradient,
4. a mobility entry in \(D_x(q)\),
5. a role in \(U_x(q)\), \(P_x(q)\), a basin \(A_i\), or a memory coordinate
\(\psi_\alpha\).

\subsection{Concrete construction from the species library}

After molecule counts (Section~\ref{sec:library}) are fixed, the abstract
coordinate \(q\) above is instantiated concretely as the collection of all
site positions:
\[
q=
\{r_{m,u}\}_{m=1}^{N_{\rm mol},\,u\in\mathcal I_{s(m)}}.
\]

Here \(m\) indexes molecules, \(s(m)\) is the species of molecule \(m\), and
\[
r_{m,u}\in\mathbb R^3.
\]

The full coordinate dimension is
\[
d=3\sum_m n_{s(m)}.
\]

The configuration space is periodic:
\[
\Omega_x=
[0,L_x)\times[0,L_y)\times[0,L_z)
\]
for each site coordinate, with unwrapped coordinates stored separately for
transport. (The reduced-coordinate case is treated formally in
Section~\ref{sec:reduced}.)

\section{Charge polarization}
\label{sec:polarization}

\subsection{Abstract definition}

Define formal-charge polarization:
\[
P_x(q)=\sum_{\alpha\in\mathcal Q} z_\alpha r_\alpha^{\rm unwrap}(q)
\in\R^3.
\]

Here:
\[
\mathcal Q
=
\text{set of charge sites},
\]
\[
z_\alpha
=
\text{dimensionless charge number of site }\alpha,
\]
\[
r_\alpha^{\rm unwrap}(q)
=
\text{unwrapped position of charge site }\alpha.
\]

The gradient of charge polarization is
\[
G_P(q)=\nabla_q P_x(q)\in\R^{3\times d}.
\]

Component notation:
\[
(G_P)_{a\ell}(q)
=
\frac{\partial P_{x,a}(q)}{\partial q_\ell},
\]
where \(a\in\{x,y,z\}\).

Using formal charge units, the final conductivity prefactor is
\[
\frac{F^2}{RT}.
\]

If physical charge units are used instead, the prefactor changes to
\[
\frac{1}{k_BT\,V}.
\]

This document uses formal charge units.

\subsection{Concrete construction from species topology}

For every molecule \(m\) and site \(u\), the charge number is
\[
z_{m,u}=z_{s(m),u}
\]
(taken from the charge-site model \(\mathcal Q_{s}\) of Section~\ref{sec:library}).

The unwrapped site position is
\[
R_{m,u}^{\rm unwrap}(q).
\]

The formal-charge polarization is then
\[
P_x(q)
=
\sum_{m}\sum_{u\in\mathcal I_{s(m)}}
z_{m,u}R_{m,u}^{\rm unwrap}(q)
\in\mathbb R^3,
\]
which instantiates the abstract sum over \(\mathcal Q\) above with
\(\mathcal Q=\{(m,u)\}\).

The gradient is
\[
G_P(q)=\nabla_qP_x(q)\in\R^{3\times d}.
\]

For coordinate component \(r_{m,u,\ell}\), where \(\ell\in\{x,y,z\}\),
\[
\frac{\partial P_{x,a}}{\partial r_{m,u,\ell}}
=
z_{m,u}\delta_{a\ell}.
\]

If a reduced coordinate \(\xi(q)\) is used instead of atomistic coordinates,
then the reduced charge-polarization gradient is obtained by the chain rule:
\[
\nabla_{\xi}P_x
=
\nabla_qP_x
\left(\frac{\partial \xi}{\partial q}\right)^{\#}.
\]

\section{Force-field / free-energy construction}
\label{sec:forcefield}

\subsection{Abstract decomposition}

Define molar free energy
\[
U_x:\Omega_x\to\R
\qquad [\mathrm{J\,mol^{-1}}].
\]

It may be decomposed as
\[
U_x(q)
=
U_{\rm bonded}
+
U_{\rm rep}
+
U_{\rm disp}
+
U_{\rm Coul}
+
U_{\rm pol}
+
U_{\rm solv}
+
U_{\rm coord}
+
U_{\rm pack}
+
U_{\rm atm}
+
U_{\rm ext}.
\]

The exact decomposition is not important in the abstract generator formalism
(Section~\ref{sec:generator}); what matters there is that \(U_x(q)\) is a
numerical function evaluable at quadrature points. The subsections below give
the concrete construction of that function from the species library
\(\mathcal L\), so that \(U_x\) is no longer a free input but a derived
quantity.

\subsection{Intramolecular bonded energy}

Bond term:
\[
U_{\rm bond}(q)
=
N_A
\sum_m
\sum_{(u,v)\in\mathcal E_{s(m)}^{\rm bond}}
\frac12 k_{uv}^{\rm bond}
\left(
|r_{m,u}-r_{m,v}|-r_{uv}^0
\right)^2.
\]

Angle term:
\[
U_{\rm angle}(q)
=
N_A
\sum_m
\sum_{(u,v,w)\in\mathcal E_{s(m)}^{\rm angle}}
\frac12 k_{uvw}^{\rm angle}
\left(
\theta_{uvw}(q)-\theta_{uvw}^0
\right)^2.
\]

Torsion term:
\[
U_{\rm torsion}(q)
=
N_A
\sum_m
\sum_{(u,v,w,z)\in\mathcal E_{s(m)}^{\rm torsion}}
\sum_{n=1}^{n_{\max}}
\frac12 V_n
\left[
1+\cos(n\varphi_{uvwz}(q)-\gamma_n)
\right].
\]

Thus:
\[
U_{\rm intra}=U_{\rm bond}+U_{\rm angle}+U_{\rm torsion}.
\]

\subsection{Lennard--Jones / short-range energy}

For two sites \(a=(m,u)\), \(b=(n,v)\), with \(a<b\), define minimum-image
separation:
\[
r_{ab}=|r_a-r_b|_{\rm PBC}.
\]

Mixing rules:
\[
\sigma_{ab}=\frac{\sigma_a+\sigma_b}{2},
\qquad
\epsilon_{ab}=\sqrt{\epsilon_a\epsilon_b}.
\]

Lennard--Jones energy:
\[
U_{\rm LJ}(q)
=
N_A
\sum_{a<b}
4\epsilon_{ab}
\left[
\left(\frac{\sigma_{ab}}{r_{ab}}\right)^{12}
-
\left(\frac{\sigma_{ab}}{r_{ab}}\right)^6
\right]f_{\rm cut}(r_{ab}).
\]

The cutoff switching function is:
\[
f_{\rm cut}(r)=
\begin{cases}
1, & r\le r_{\rm on},\\
1-10s^3+15s^4-6s^5, & r_{\rm on}<r<r_{\rm off},\\
0, & r\ge r_{\rm off},
\end{cases}
\]
where
\[
s=\frac{r-r_{\rm on}}{r_{\rm off}-r_{\rm on}}.
\]

\subsection{Electrostatics}

Let
\[
q_a^{\rm SI}=e z_a.
\]

For periodic boundaries, use Ewald electrostatics:
\[
U_{\rm Coul}(q)
=
N_A
\frac{1}{4\pi\epsilon_0\epsilon_\infty}
\left[
U_{\rm real}
+
U_{\rm reciprocal}
+
U_{\rm self}
+
U_{\rm surface}
\right].
\]

Real-space term:
\[
U_{\rm real}
=
\frac12
\sum_{a\ne b}
q_a^{\rm SI}q_b^{\rm SI}
\frac{\operatorname{erfc}(\alpha r_{ab})}{r_{ab}}.
\]

Reciprocal-space term:
\[
U_{\rm reciprocal}
=
\frac{2\pi}{V_{\rm cell}}
\sum_{k\ne0}
\frac{\exp[-k^2/(4\alpha^2)]}{k^2}
\left|
\sum_a q_a^{\rm SI}\exp(i k\cdot r_a)
\right|^2.
\]

Self term:
\[
U_{\rm self}
=
-\frac{\alpha}{\sqrt\pi}
\sum_a (q_a^{\rm SI})^2.
\]

For conducting boundary conditions:
\[
U_{\rm surface}=0.
\]

For vacuum boundary condition:
\[
U_{\rm surface}
=
\frac{2\pi}{3V_{\rm cell}}
\left|
\sum_a q_a^{\rm SI}r_a
\right|^2.
\]

\subsection{Polarization, if used}

Let induced dipoles be \(\mu_a\), polarizabilities \(\alpha_a\), and electric
fields \(E_a(q)\). The polarization energy is
\[
U_{\rm pol}(q,\{\mu_a\})
=
N_A
\left[
\sum_a
\frac12
\mu_a^\top\alpha_a^{-1}\mu_a
-
\sum_a \mu_a^\top E_a(q)
+
\frac12
\sum_{a\ne b}\mu_a^\top T_{ab}\mu_b
\right].
\]

The induced dipoles are minimized out:
\[
\{\mu_a^\star(q)\}
=
\argmin_{\{\mu_a\}}U_{\rm pol}(q,\{\mu_a\}).
\]

Then:
\[
U_{\rm pol}^{\rm eff}(q)
=
U_{\rm pol}(q,\{\mu_a^\star(q)\}).
\]

If polarization is not used:
\[
U_{\rm pol}^{\rm eff}(q)=0.
\]

\subsection{Solvation and coordination free energy}

Define smooth coordination number between center \(a\) and ligand class \(L\):
\[
n_{aL}(q)
=
\sum_{b\in L}
s_{aL}(r_{ab}),
\]
where
\[
s_{aL}(r)
=
\frac{1}{1+(r/r_{aL}^0)^{p_{aL}}}.
\]

Coordination energy:
\[
U_{\rm coord}(q)
=
N_A
\sum_{a,L}
\Delta g_{aL}^{\rm coord}
\,n_{aL}(q).
\]

Solvation-shell competition energy:
\[
U_{\rm solv}(q)
=
N_A
\sum_a
F_a^{\rm solv}
\left(
n_{aL_1}(q),
n_{aL_2}(q),
\ldots
\right).
\]

A concrete convex shell penalty is:
\[
F_a^{\rm solv}(n_1,\ldots,n_m)
=
\frac12
\sum_{\ell,\ell'}
K^a_{\ell\ell'}
(n_\ell-\bar n_\ell^a)
(n_{\ell'}-\bar n_{\ell'}^a).
\]

\subsection{Packing / excluded-volume free energy}

Let local occupied volume fraction be
\[
\varphi(q)
=
\sum_a v_a W_a(q),
\]
where \(v_a\) is molecular volume and \(W_a(q)\) is a normalized local kernel.

Carnahan--Starling-like packing penalty:
\[
U_{\rm pack}(q)
=
N_A RT
\sum_{\ell}
w_\ell^{\rm grid}
\frac{
\varphi_\ell(4-3\varphi_\ell)
}{
(1-\varphi_\ell)^2
}.
\]

\subsection{Total effective free energy}

The final free energy is:
\[
U_x(q)
=
U_{\rm intra}
+
U_{\rm LJ}
+
U_{\rm Coul}
+
U_{\rm pol}^{\rm eff}
+
U_{\rm coord}
+
U_{\rm solv}
+
U_{\rm pack}.
\]

This is the concrete realization of the abstract decomposition given at the
start of this section (\(U_{\rm bonded}\to U_{\rm intra}\),
\(U_{\rm rep}+U_{\rm disp}\to U_{\rm LJ}\), \(U_{\rm atm}\) and
\(U_{\rm ext}\) absorbed into \(U_{\rm coord}/U_{\rm solv}\) as needed).

The equilibrium measure is:
\[
\mu_x(\dd q)
=
Z_x^{-1}\exp[-U_x(q)/(RT)]\,\dd q.
\]

\section{Mobility tensor construction}
\label{sec:mobility}

\subsection{Abstract requirement}

Define
\[
D_x(q)\in\R^{d\times d}
\qquad [\mathrm{m^2\,s^{-1}}].
\]

It must satisfy
\[
D_x(q)=D_x(q)^\top,
\qquad
D_x(q)\succeq0.
\]

If built from resistance:
\[
D_x(q)=k_BT\,R_x(q)^\#.
\]

Resistance may be assembled abstractly as
\[
R_x(q)
=
R_{\rm hydro}
+
R_{\rm crowd}
+
R_{\rm charge\ cloud}
+
R_{\rm atmosphere}
+
R_{\rm cage}
+
R_{\rm constraint}.
\]

The subsections below give the concrete construction of \(R_x(q)\) (and hence
\(D_x(q)\)) from the mobility/friction library \(\mathcal M\subset\mathcal L\).

\subsection{Resistance tensor}

Define a generalized resistance matrix
\[
R_x(q)\in\mathbb R^{d\times d},
\]
decomposed concretely as:
\[
R_x(q)
=
R_{\rm Stokes}(q)
+
R_{\rm hydro}(q)
+
R_{\rm constraint}(q)
+
R_{\rm atmosphere}(q)
+
R_{\rm cage}(q)
+
R_{\rm microvisc}(q).
\]

\subsection{Local viscosity}

Let the local viscosity be:
\[
\eta(q)
=
\eta_{\rm bulk}(x)
\,
f_{\rm salt}(I(q))
\,
f_{\rm additive}(q)
\,
f_{\rm free\ volume}(\varphi(q)).
\]

A concrete form:
\[
f_{\rm salt}(I)=1+A_{\eta}\sqrt{I}+B_{\eta}I,
\]
\[
f_{\rm additive}(q)
=
1+
\sum_L a_L n_L(q),
\]
\[
f_{\rm free\ volume}(\varphi)
=
\left(1-\frac{\varphi}{\varphi_{\max}}\right)^{-\gamma}.
\]

\subsection{Stokes resistance}

For site \(a\):
\[
\zeta_a(q)=6\pi\eta(q)a_a.
\]

The diagonal Stokes resistance is:
\[
R_{\rm Stokes,ab}(q)
=
\zeta_a(q)I_3\delta_{ab}.
\]

\subsection{Hydrodynamic cross mobility}

If hydrodynamic interactions are included, use Rotne--Prager--Yamakawa mobility
\[
H_{ab}(q)
=
\begin{cases}
\frac{1}{6\pi\eta a_a}I_3, & a=b,\\[4pt]
\frac{1}{8\pi\eta r_{ab}}
\left[
I_3+\hat r_{ab}\hat r_{ab}^{\top}
+
\frac{2(a_a^2+a_b^2)}{3r_{ab}^2}
\left(
I_3-3\hat r_{ab}\hat r_{ab}^{\top}
\right)
\right], & a\ne b,\ r_{ab}>a_a+a_b.
\end{cases}
\]

Then
\[
D_{\rm hydro}(q)=k_BT\,H(q).
\]

If using resistance form:
\[
R_{\rm hydro}(q)=k_BT\,D_{\rm hydro}(q)^\#.
\]

\subsection{Constraint resistance}

For a constraint mode \(m\) with direction vector \(s_m(q)\), lifetime
\(\tau_m(q)\), and length \(\ell_m(q)\), define:
\[
R_{\rm constraint}^{(m)}(q)
=
k_BT
\frac{\tau_m(q)}{\ell_m(q)^2}
s_m(q)s_m(q)^\top.
\]

Sum:
\[
R_{\rm constraint}(q)=\sum_m R_{\rm constraint}^{(m)}(q).
\]

\subsection{Atmosphere resistance}

Let charge density mode be
\[
\rho_z(k;q)
=
\sum_a z_a \exp(i k\cdot r_a).
\]

For Debye screening:
\[
\kappa^2(q)
=
\frac{F^2}{\epsilon_0\epsilon_r RT}
\sum_i z_i^2 c_i(q).
\]

Atmosphere resistance projected onto charge form factors:
\[
R_{\rm atmosphere}(q)
=
\sum_{k\in\mathcal K}
\Gamma(k;\kappa,\eta,\epsilon_r)
\,g_k(q)g_k(q)^\top,
\]
where \(g_k(q)=\nabla_q\rho_z(k;q)\) and \(\Gamma\ge0\) is the finite-size
PNP--Stokes response kernel.

A simple admissible diagonal kernel is:
\[
\Gamma(k;\kappa,\eta,\epsilon_r)
=
\Gamma_0
\frac{k^2}{(k^2+\kappa^2)^2}
\eta.
\]

\subsection{Total mobility}

The diffusion/mobility tensor is:
\[
D_x(q)=k_BT\,R_x(q)^\#.
\]

Validation:
\[
D_x(q)=D_x(q)^\top,\qquad D_x(q)\succeq0.
\]

\section{Generator}
\label{sec:generator}

With \(U_x\) (Section~\ref{sec:forcefield}) and \(D_x\)
(Section~\ref{sec:mobility}) now concretely constructed from the species
library \(\mathcal L\), the reversible generator is well defined.

\subsection{Thermal factor}

Use the molar gas constant
\[
R=8.31446261815324\ \mathrm{J\,mol^{-1}\,K^{-1}}.
\]

Define
\[
\beta_{\rm mol}=\frac{1}{RT}.
\]

Then
\[
\exp[-\beta_{\rm mol}U_x(q)]
\]
is the Boltzmann factor.

\subsection{Reversible generator}

For overdamped Smoluchowski dynamics:
\[
L_x f(q)
=
\frac{1}{\mu_x(q)}
\nabla_q\cdot
\left(
\mu_x(q)D_x(q)\nabla_q f(q)
\right).
\]

Equivalently:
\[
L_x f
=
D_x:\nabla_q\nabla_q f
+
\left[
\nabla_q\cdot D_x
-
\beta_{\rm mol}D_x\nabla_q U_x
\right]\cdot\nabla_q f.
\]

The equilibrium measure is
\[
\mu_x(\dd q)
=
Z_x^{-1}
\exp[-\beta_{\rm mol}U_x(q)]
\dd q.
\]

The model does not need to simulate \(L_x\) in production. It uses \(L_x\)
through quadrature, committor/capacity solves, transition-path moment solves,
and Mori inner products.

\section{Reduced generator derived from full atomistic generator}
\label{sec:reduced}

If the model uses reduced coordinates
\[
\xi=\Xi(q)\in\mathbb R^m,
\]
then the reduced PMF is not arbitrary. It is:
\[
U_{\rm red}(\xi)
=
-RT
\log
\int_{\Omega_x}
\delta(\Xi(q)-\xi)
\exp[-U_x(q)/(RT)]\,\dd q.
\]

The reduced mobility is the conditional projection:
\[
D_{\rm red}(\xi)
=
\E_{\mu_x}
\left[
\nabla_q\Xi(q)D_x(q)\nabla_q\Xi(q)^\top
\mid
\Xi(q)=\xi
\right].
\]

The reduced polarization is:
\[
P_{\rm red}(\xi)
=
\E_{\mu_x}
\left[
P_x(q)\mid\Xi(q)=\xi
\right].
\]

The reduced basins are:
\[
A_i^{\rm red}
=
\{\xi:\xi\in B_i\},
\]
where
\[
B_i=\Xi(A_i).
\]

Thus, a reduced model is first-principles only if
\[
U_{\rm red},D_{\rm red},P_{\rm red},A_i^{\rm red},\psi^{\rm red}
\]
are derived by conditional projection from the atomistic generator, not fitted
as independent descriptor functions. Every reduced-coordinate object used
elsewhere in this document (Section~\ref{sec:configspace}) must satisfy this
requirement.

\section{Projection basins}
\label{sec:basins}

Choose \(N\) disjoint basins:
\[
\mathcal A=\{A_1,\ldots,A_N\},
\]
where
\[
A_i\subset\Omega_x,
\qquad
A_i\cap A_j=\varnothing \quad(i\ne j).
\]

Every coordinate \(q\) used in the finite model is assigned to exactly one
basin:
\[
s(q)=i
\quad
\Longleftrightarrow
\quad
q\in A_i.
\]

Each basin \(A_i\) has a stoichiometry vector:
\[
\nu_i=(\nu_{i,a})_{a\in\mathcal C}.
\]

\[
\nu_{i,a}
=
\text{number of conserved component }a\text{ in basin/motif }i.
\]

Examples:
\[
\nu_{i,\mathrm{Li}}=1
\]
for a motif containing one Li.
\[
\nu_{i,\mathrm{PF6}}=1
\]
for a motif containing one PF6.
\[
\nu_{i,\mathrm{FEC}}=1
\]
for a Li--FEC ligand motif.

A basin is a coordinate set, not a terminal correction.

\subsection{Concrete basin definitions from physical coordinates}

\subsubsection{Switching functions}

Define smooth pair contact:
\[
s_{ab}(q)
=
\frac{1}{1+(r_{ab}(q)/r_{ab}^0)^p}.
\]

Define coordination number:
\[
n_{aL}(q)=\sum_{b\in L}s_{ab}(q).
\]

Define ion-pair graph edge:
\[
E_{a b}^{\rm pair}(q)
=
\mathbf 1\{r_{ab}(q)<r_{\rm pair}^{\rm cut}\}.
\]

Define ligand-shell indicator:
\[
E_{aL}^{\rm lig}(q)
=
\mathbf 1\{n_{aL}(q)\ge n_{aL}^{\rm cut}\}.
\]

Define cluster graph:
\[
G(q)=(V,E(q)),
\]
where vertices are charged species and
\[
(a,b)\in E(q)
\iff
E_{ab}^{\rm pair}(q)=1.
\]

Cluster membership is connected component membership in \(G(q)\).

\subsubsection{Pair basins}

For Li--anion pair \((\ell,a)\), define:
\[
r_{\ell a}=|r_\ell-r_a|.
\]

Contact ion pair:
\[
A_{\rm CIP}
=
\{q:r_{\ell a}<r_{\rm CIP}\}.
\]

Solvent-separated ion pair:
\[
A_{\rm SSIP}
=
\{q:r_{\rm CIP}\le r_{\ell a}<r_{\rm SSIP}\}.
\]

Free state:
\[
A_{\rm free}
=
\{q:r_{\ell a}\ge r_{\rm SSIP}\ \text{for all nearby anions}\}.
\]

\subsubsection{Ligand-conditioned basins}

Neutral ligand occupancy around Li:
\[
n_{\ell L}(q)=\sum_{b\in L}s_{\ell b}(q).
\]

Li--ligand basin:
\[
A_{\rm LiL}
=
\{q:n_{\ell L}(q)\ge n_{\rm lig}^{\rm cut}\}.
\]

Ligand-separated SSIP:
\[
A_{\rm LiL\cdots A}
=
A_{\rm SSIP}
\cap
A_{\rm LiL}
\cap
\{r_{\ell a}\in[r_{\rm CIP},r_{\rm SSIP})\}.
\]

\subsubsection{Anion-orientation basins}

Let \(u_a(q)\) be the principal axis of anion \(a\), and
\[
\hat r_{\ell a}=\frac{r_a-r_\ell}{|r_a-r_\ell|}.
\]

Define orientation coordinate:
\[
\omega_a(q)=u_a(q)\cdot \hat r_{\ell a}.
\]

Orientation basin:
\[
A_{\omega,k}
=
\{q:\omega_k^-\le\omega_a(q)<\omega_k^+\}.
\]

\subsubsection{Final basin index}

The final basin is the joint bin:
\[
A_i=
A_{\rm pair}
\cap
A_{\rm ligand}
\cap
A_{\rm cluster}
\cap
A_{\rm orientation}
\cap
A_{\rm cage}
\cap
A_{\rm atmosphere}.
\]

The set of all nonempty joint bins forms:
\[
\mathcal A=\{A_1,\ldots,A_N\}.
\]

This concrete construction instantiates the abstract disjoint decomposition
\(\mathcal A=\{A_1,\ldots,A_N\}\) given at the start of this section, together
with the stoichiometry vectors \(\nu_i\).

\section{Quadrature}
\label{sec:quadrature}

For each basin \(A_i\), choose quadrature points and weights:
\[
\{(q_{i\ell},w_{i\ell})\}_{\ell=1}^{n_i}.
\]

Each
\[
q_{i\ell}\in A_i.
\]

Weights satisfy
\[
\int_{A_i} f(q)\dd q
\approx
\sum_{\ell=1}^{n_i}w_{i\ell}f(q_{i\ell}).
\]

Weights may be generated by:

1. tensor quadrature,
2. sparse grid quadrature,
3. finite element quadrature,
4. Monte Carlo / importance quadrature,
5. umbrella-sampling reweighting.

The model only requires that the same weights are used consistently in all
basin integrals.

\section{Mass balance through chemical potentials}
\label{sec:massbalance}

\subsection{Why chemical potentials are needed}

Composition gives conserved totals \(C_a\). The projected basins \(A_i\) are
motif states that consume conserved components according to \(\nu_{i,a}\).
Therefore the basin populations \(c_i\) must satisfy:
\[
\sum_i\nu_{i,a}c_i=C_a
\qquad
\text{for every conserved component }a.
\]

This is achieved by component chemical potentials \(\lambda_a\). (This is the
same requirement obtained if one starts directly from the composition
\(x=(r,\mathcal L)\) and asks for basin populations consistent with
conservation of \(\{C_a\}\); the construction below is the unique consistent
solution.)

\subsection{Grand-canonical restricted basin weight}

For each basin \(i\), define
\[
Z_i(\lambda)
=
\int_{A_i}
\exp
\left[
-\beta_{\rm mol}U_x(q)
+
\sum_{a\in\mathcal C}\nu_{i,a}\lambda_a
\right]\dd q.
\]

The dimensionless chemical potential \(\lambda_a\) is measured in units of
\(RT\). It is not in J/mol. If using molar chemical potential
\(\mu_a^{\rm chem}\), then
\[
\lambda_a=\frac{\mu_a^{\rm chem}}{RT}.
\]

Use standard concentration
\[
C^\circ=1000\ \mathrm{mol\,m^{-3}}.
\]

Then basin concentration is
\[
c_i(\lambda)=C^\circ Z_i(\lambda).
\]

Quadrature form:
\[
Z_i(\lambda)
\approx
\sum_{\ell=1}^{n_i}
w_{i\ell}
\exp
\left[
-\beta_{\rm mol}U_x(q_{i\ell})
+
\sum_{a\in\mathcal C}\nu_{i,a}\lambda_a
\right].
\]

\subsection{Mass-balance residual}

For every conserved component \(a\), define
\[
r_a(\lambda)
=
\sum_{i=1}^N\nu_{i,a}c_i(\lambda)-C_a.
\]

The chemical potential solve is:
\[
r_a(\lambda^\star)=0
\qquad
\forall a.
\]

\subsection{Jacobian}

The derivative is
\[
\frac{\partial c_i}{\partial\lambda_b}
=
\nu_{i,b}c_i.
\]

Therefore
\[
J_{ab}(\lambda)
=
\frac{\partial r_a}{\partial\lambda_b}
=
\sum_{i=1}^N\nu_{i,a}\nu_{i,b}c_i(\lambda).
\]

\subsection{Newton algorithm}

Input:
\[
C_a,\quad \nu_{i,a},\quad U_x,\quad A_i,\quad q_{i\ell},\quad w_{i\ell}.
\]

Tolerance:
\[
\tau_{\rm mass}=10^{-10}.
\]

Maximum iterations:
\[
N_{\rm Newton}=100.
\]

Initialization:
\[
\lambda_a^{(0)}
=
\log
\left(
\frac{\max(C_a,C_{\min})}{C^\circ}
\right),
\]
with
\[
C_{\min}=10^{-30}\ \mathrm{mol\,m^{-3}}.
\]

Iteration \(k\):

1. Compute all \(Z_i(\lambda^{(k)})\).

2. Compute all
\[
c_i^{(k)}=C^\circ Z_i(\lambda^{(k)}).
\]

3. Compute residual
\[
r_a^{(k)}
=
\sum_i\nu_{i,a}c_i^{(k)}-C_a.
\]

4. Compute normalized residual
\[
\eta^{(k)}
=
\max_a
\frac{|r_a^{(k)}|}
{\max(C_a,C_{\min})}.
\]

5. Stop if
\[
\eta^{(k)}<\tau_{\rm mass}.
\]

6. Compute Jacobian
\[
J_{ab}^{(k)}
=
\sum_i\nu_{i,a}\nu_{i,b}c_i^{(k)}.
\]

7. Solve the linear Newton system
\[
J^{(k)}\Delta\lambda^{(k)}=-r^{(k)}.
\]

If \(J^{(k)}\) is singular or ill-conditioned, use Tikhonov-regularized solve:
\[
\left(J^{(k)}+\epsilon_J I\right)\Delta\lambda^{(k)}=-r^{(k)},
\]
where
\[
\epsilon_J=10^{-12}\max(1,\|J^{(k)}\|_2).
\]

8. Backtracking line search:

Start with
\[
\alpha=1.
\]

Trial:
\[
\lambda_{\rm trial}=\lambda^{(k)}+\alpha\Delta\lambda^{(k)}.
\]

Compute
\[
\eta_{\rm trial}.
\]

Accept the step if
\[
\eta_{\rm trial}<\eta^{(k)}.
\]

If not accepted, set
\[
\alpha\leftarrow \frac12\alpha
\]
and try again.

Fail if
\[
\alpha<2^{-40}.
\]

9. Update:
\[
\lambda^{(k+1)}=\lambda_{\rm trial}.
\]

After convergence:
\[
\lambda^\star=\lambda^{(k)}.
\]

The final concentration is:
\[
c_i=c_i(\lambda^\star).
\]

\subsection{Density quadrature weights}

Define the concentration-density quadrature weight:
\[
W_{i\ell}
=
C^\circ
w_{i\ell}
\exp
\left[
-\beta_{\rm mol}U_x(q_{i\ell})
+
\sum_a\nu_{i,a}\lambda_a^\star
\right].
\]

Then
\[
c_i=\sum_{\ell=1}^{n_i}W_{i\ell}.
\]

These \(W_{i\ell}\) are used in all basin averages.

\section{Reactive fluxes}
\label{sec:fluxes}

\subsection{Purpose}

Reactive fluxes determine how often the projected process moves between basins.
They produce:
\[
K_{ij},\quad Q_{ij}.
\]
\[
K_{ij}
\quad [\mathrm{mol\,m^{-3}\,s^{-1}}]
\]
is the symmetric concentration flux between basin \(i\) and basin \(j\).
\[
Q_{ij}
\quad [\mathrm{s^{-1}}]
\]
is the finite-state transition rate from \(i\) to \(j\).

\subsection{Committor definition}

For an edge \((i,j)\), define the committor:
\[
q_{ij}(q)
=
\Pr_q(\tau_j<\tau_i).
\]

Here
\[
\tau_i=\inf\{t:X_t\in A_i\}.
\]

The committor satisfies:
\[
L_x q_{ij}=0
\quad
\text{on }\Omega_{ij}=\Omega_x\setminus(A_i\cup A_j),
\]
\[
q_{ij}=0
\quad
\text{on }A_i,
\]
\[
q_{ij}=1
\quad
\text{on }A_j.
\]

\subsection{How to solve the committor numerically}

Choose finite element basis functions
\[
\{\varphi_m(q)\}_{m=1}^{M}
\]
on the transition domain \(\Omega_{ij}\).

Approximate:
\[
q_{ij}(q)\approx \sum_{m=1}^{M}u_m\varphi_m(q).
\]

Boundary nodes on \(A_i\) are fixed to \(0\).
Boundary nodes on \(A_j\) are fixed to \(1\).

For interior degrees of freedom, solve:
\[
S u = f,
\]
where the stiffness matrix is
\[
S_{mn}
=
\int_{\Omega_{ij}}
\nabla\varphi_m(q)^\top
D_x(q)
\nabla\varphi_n(q)
\,\rho_x(\dd q),
\]
and \(\rho_x\) is the concentration measure:
\[
\rho_x(\dd q)
=
C^\circ
\exp
\left[
-\beta_{\rm mol}U_x(q)
+
\sum_a\nu_{s(q),a}\lambda_a^\star
\right]\dd q.
\]

The right-hand side \(f\) comes only from known boundary values after
eliminating Dirichlet nodes.

After solving \(u\), evaluate
\[
\nabla q_{ij}(q)
\]
at transition quadrature points.

\subsection{Capacity flux}

The symmetric capacity flux is:
\[
K_{ij}
=
\int_{\Omega_{ij}}
\nabla q_{ij}(q)^\top
D_x(q)
\nabla q_{ij}(q)\,
\rho_x(\dd q).
\]

Using transition quadrature points \(q_{ij,\ell}^{\Sigma}\), weights
\(w_{ij,\ell}^{\Sigma}\), and committor gradients
\[
g_{ij,\ell}=\nabla q_{ij}(q_{ij,\ell}^{\Sigma}),
\]
compute:
\[
K_{ij}
\approx
\sum_{\ell}
W_{ij,\ell}^{\Sigma}
g_{ij,\ell}^{\top}
D_x(q_{ij,\ell}^{\Sigma})
g_{ij,\ell}.
\]

The transition-surface concentration weight is:
\[
W_{ij,\ell}^{\Sigma}
=
C^\circ
w_{ij,\ell}^{\Sigma}
\exp
\left[
-\beta_{\rm mol}U_x(q_{ij,\ell}^{\Sigma})
+
\sum_a\nu_{s(q_{ij,\ell}^{\Sigma}),a}\lambda_a^\star
\right].
\]

Set:
\[
K_{ji}=K_{ij}.
\]

Set:
\[
K_{ii}=0.
\]

\subsection{Finite generator}

For \(i\ne j\):
\[
Q_{ij}=\frac{K_{ij}}{c_i}.
\]

For diagonal entries:
\[
Q_{ii}=-\sum_{j\ne i}Q_{ij}.
\]

Validation:

1. Row sums:
\[
\sum_j Q_{ij}=0.
\]

2. Detailed balance:
\[
c_iQ_{ij}=c_jQ_{ji}.
\]

3. Stationarity:
\[
\sum_i c_iQ_{ij}=0.
\]

\section{Transition-path displacement moments}
\label{sec:tpmoments}

\subsection{Purpose}

Transition-path moments determine charge transported during state changes.
They produce:
\[
d_{ij},\quad M_{ij}.
\]
\[
d_{ij}\in\R^3
\quad [\mathrm{m}]
\]
is mean formal-charge displacement for transition \(i\to j\).
\[
M_{ij}\in\R^{3\times3}
\quad [\mathrm{m^2}]
\]
is second moment of formal-charge displacement for transition \(i\to j\).

\subsection{Doob-conditioned drift and reactive path sampling}

For edge \((i,j)\), use the committor \(q_{ij}\).

The process conditioned to reach \(A_j\) before \(A_i\) has Doob-transformed
drift:
\[
b^{(ij)}(q)
=
b(q)+2D_x(q)\nabla\log q_{ij}(q),
\]
where
\[
b(q)=\nabla_q\cdot D_x(q)-\beta_{\rm mol}D_x(q)\nabla_q U_x(q)
\]
is the original overdamped drift.

The conditioned short-path SDE is:
\[
\dd q_t
=
b^{(ij)}(q_t)\dd t
+
\sqrt{2D_x(q_t)}\dd W_t.
\]

Initial points are sampled from the reactive exit distribution on the
\(i\)-side boundary:
\[
p_{\rm exit}(q)
\propto
\rho_x(q)
\,n_{ij}(q)^\top D_x(q)\nabla q_{ij}(q),
\]
where \(n_{ij}\) is the outward normal from \(A_i\) toward \(A_j\).

For each sampled path:

1. Start at \(q_0\) from \(p_{\rm exit}\).

2. Integrate the conditioned SDE (Euler--Maruyama)
\[
q_{n+1}
=
q_n
+
b^{(ij)}(q_n)\Delta t
+
\sqrt{2D_x(q_n)\Delta t}\,\eta_n,
\qquad
\eta_n\sim\mathcal N(0,I),
\]
until hitting \(A_j\), i.e.\ until
\[
\tau_j=\inf\{t:q_t\in A_j\}.
\]

3. Record endpoint \(q_\tau\).

4. Compute charge displacement:
\[
\Delta P_{ij,\ell}=P_x(q_\tau)-P_x(q_0).
\]

This is a local transition-path calculation, not a long-time GK integral.

\subsection{Weighted moment estimator}

Given path samples
\[
\{(\Delta P_{ij,\ell},w_{ij,\ell}^{\rm path})\}_{\ell=1}^{n_{ij}^{\rm path}},
\]
compute:
\[
S_{ij}^{\rm path}=\sum_\ell w_{ij,\ell}^{\rm path}.
\]

Then:
\[
d_{ij}
=
\frac{1}{S_{ij}^{\rm path}}
\sum_\ell
w_{ij,\ell}^{\rm path}\Delta P_{ij,\ell}.
\]

\[
M_{ij}
=
\frac{1}{S_{ij}^{\rm path}}
\sum_\ell
w_{ij,\ell}^{\rm path}
\Delta P_{ij,\ell}\Delta P_{ij,\ell}^{\top}.
\]

Set reverse moments:
\[
d_{ji}=-d_{ij},
\qquad
M_{ji}=M_{ij}.
\]

Set:
\[
d_{ii}=0,
\qquad
M_{ii}=0.
\]

\subsection{Displacement-scale validation}

For every transition displacement sample, require:
\[
\|\Delta P_{ij,\ell}\| \le L_{\max}.
\]

For molecular local transitions, choose for example:
\[
L_{\max}=10^{-8}\ \mathrm{m}.
\]

A transition displacement of order meters must fail, because \(P_x\) is an
unwrapped molecular charge displacement in meters.

\section{Within-state self-current}
\label{sec:selfcurrent}

\subsection{Purpose}

Self-current captures continuous unbounded charge mobility while the projected
state does not change. It produces:
\[
D_i^{\rm self}\in\R^{3\times3}
\qquad [\mathrm{m^2\,s^{-1}}].
\]

\subsection{Unbounded transport projector}

For each basin \(A_i\), define:
\[
\Pi_i(q)\in\R^{d\times d}.
\]

\(\Pi_i(q)\) projects coordinate motion onto directions that produce
unbounded DC charge transport.

Examples:

1. Free ion translation:
\[
\Pi_i=I
\]
on that ion's translational coordinates.

2. Bound internal pair vibration:
\[
\Pi_i=0
\]
for the bounded internal coordinate.

3. Neutral pair center-of-mass translation:
\[
\Pi_i=I_{\rm COM}
\]
but
\[
G_P\Pi_i=0
\]
if the neutral pair is exactly charge-neutral and co-moving.

4. Partner-switch or identity-diffusion coordinate:
\[
\Pi_i
\]
is nonzero only along the unbounded identity coordinate.

\subsection{Self-current integral}

For each state:
\[
D_i^{\rm self}
=
\frac{1}{c_i}
\int_{A_i}
G_P(q)
\Pi_i(q)
D_x(q)
\Pi_i(q)^\top
G_P(q)^\top
\,\rho_x(\dd q).
\]

Quadrature:
\[
D_i^{\rm self}
\approx
\frac{1}{c_i}
\sum_{\ell=1}^{n_i}
W_{i\ell}
G_P(q_{i\ell})
\Pi_i(q_{i\ell})
D_x(q_{i\ell})
\Pi_i(q_{i\ell})^\top
G_P(q_{i\ell})^\top.
\]

Validation:
\[
D_i^{\rm self}=D_i^{\rm self\top},
\qquad
D_i^{\rm self}\succeq0.
\]

\subsection{Multi-center charge covariance}

If a basin contains multiple charged centers, do not use
\[
z_{\rm net}^2D_{\rm COM}.
\]

Let
\[
z_i=(z_1,\ldots,z_m)^\top
\]
be the charge vector of centers in basin \(i\), and let
\[
D_i^{\rm centers}\in\R^{m\times m}
\]
be the state-local mobility covariance of those centers along one Cartesian
axis.

Then charge mobility is:
\[
D_Q(i)=z_i^\top D_i^{\rm centers}z_i.
\]

For a Li--anion pair:
\[
z=(+1,-1)^\top.
\]

Thus:
\[
D_Q
=
D_{\rm Li}
+
D_{\rm A}
-
2D_{\rm Li,A}.
\]

If
\[
D_{\rm Li,A}<0,
\]
then the cross term raises conductivity.

\section{Direct projected diffusivity-density}
\label{sec:directB}

Define:
\[
B
=
\sum_i c_iD_i^{\rm self}
+
\frac{1}{2}\sum_i\sum_jK_{ij}M_{ij}.
\]

Units:
\[
[c_iD_i^{\rm self}]
=
\mathrm{mol\,m^{-3}}\cdot\mathrm{m^2\,s^{-1}}
=
\mathrm{mol\,m^{-1}\,s^{-1}}.
\]

\[
[K_{ij}M_{ij}]
=
\mathrm{mol\,m^{-3}\,s^{-1}}\cdot \mathrm{m^2}
=
\mathrm{mol\,m^{-1}\,s^{-1}}.
\]

So:
\[
[B]=\mathrm{mol\,m^{-1}\,s^{-1}}.
\]

\section{Finite-state memory correction}
\label{sec:finitestate}

\subsection{State drift}

For each state:
\[
b_i=\sum_jQ_{ij}d_{ij}\in\R^3.
\]

Component:
\[
b_{i,a}=\sum_jQ_{ij}d_{ij,a}.
\]

\subsection{Poisson equation}

For each Cartesian axis \(a\), solve:
\[
-Q\chi_a=b_a.
\]

Because \(Q\) has a zero eigenvalue, impose gauge:
\[
\sum_i c_i\chi_{i,a}=0.
\]

\subsection{Numerical solve}

Construct block matrix:
\[
\mathcal M
=
\begin{bmatrix}
-Q & \one\\
c^\top & 0
\end{bmatrix}.
\]

Construct right-hand side:
\[
y_a=
\begin{bmatrix}
b_a\\
0
\end{bmatrix}.
\]

Solve:
\[
\mathcal M
\begin{bmatrix}
\chi_a\\
\lambda_a^{Q}
\end{bmatrix}
=
y_a.
\]

Here \(\lambda_a^{Q}\) is a Lagrange multiplier enforcing the gauge.

The solution vector is:
\[
\chi_a=(\chi_{1,a},\ldots,\chi_{N,a})^\top.
\]

\subsection{Correction tensor}

\[
C^Q_{ab}
=
\sum_i c_i b_{i,a}\chi_{i,b}.
\]

Thus:
\[
C^Q\in\R^{3\times3}.
\]

Symmetrize numerically:
\[
C^Q\leftarrow\frac12(C^Q+C^{Q\top}).
\]

\section{Continuous Mori memory correction}
\label{sec:mori}

\subsection{Memory coordinates}

Choose memory coordinates:
\[
\Psi(q)
=
(\psi_1(q),\ldots,\psi_m(q))^\top.
\]

Examples:
\[
\psi_{\rm atm}(q)
=
\text{ionic-atmosphere polarization coordinate},
\]
\[
\psi_{\rm cage}(q)
=
\text{solvent-cage/backjump coordinate},
\]
\[
\psi_{\rm partner}(q)
=
\text{partner residence coordinate},
\]
\[
\psi_{\rm orient}(q)
=
\text{anion orientation coordinate}.
\]

\subsection{Concrete memory coordinate definitions}

\subsubsection{Charge-density / atmosphere modes}

For wavevector \(k\):
\[
\psi_{k,c}(q)=
\sum_a z_a\cos(k\cdot r_a),
\]
\[
\psi_{k,s}(q)=
\sum_a z_a\sin(k\cdot r_a).
\]

\subsubsection{Cage coordinate}

Let \(r_\ell\) be Li position and \(\mathcal N_\ell(q)\) its solvent cage.
Cage center:
\[
R_{\rm cage}(q)
=
\frac{\sum_{b\in\mathcal N_\ell(q)}w_b r_b}
{\sum_{b\in\mathcal N_\ell(q)}w_b}.
\]

Cage displacement coordinate:
\[
\psi_{\rm cage}(q)=
(r_\ell-R_{\rm cage})\cdot e_{\rm cage},
\]
where \(e_{\rm cage}\) is a declared cage-axis basis vector.

\subsubsection{Partner-residence coordinate}

For current partner \(a^\star(q)\):
\[
\psi_{\rm partner}(q)
=
s_{\ell a^\star}(q).
\]

\subsubsection{Ligand-shell residence coordinate}

\[
\psi_{\rm ligand}(q)
=
n_{\ell L}(q)-\bar n_{\ell L}.
\]

\subsubsection{Anion orientation memory}

\[
\psi_{\rm orient}(q)
=
u_a(q)\cdot \hat r_{\ell a}.
\]

\subsubsection{Local free-volume memory}

\[
\psi_{\rm fv}(q)
=
\varphi(q)-\bar\varphi.
\]

Collect:
\[
\Psi(q)
=
(\psi_1(q),\ldots,\psi_m(q))^\top.
\]

\subsection{State-indicator orthogonalization}

For each state \(A_i\):
\[
\bar\psi_i
=
\frac{1}{c_i}
\int_{A_i}\Psi(q)\rho_x(\dd q).
\]

Quadrature:
\[
\bar\psi_i
=
\frac{1}{c_i}
\sum_{\ell}W_{i\ell}\Psi(q_{i\ell}).
\]

Define:
\[
\Psi^\perp(q)=\Psi(q)-\bar\psi_{s(q)}.
\]

This prevents double-counting memory already represented by finite state
indicators.

If \(\bar\psi_{s(q)}\) is piecewise constant, then inside a basin:
\[
\nabla_q\Psi^\perp(q)=\nabla_q\Psi(q).
\]

At basin boundaries, discontinuities are already handled by the finite-state
Poisson correction (Section~\ref{sec:finitestate}).

\subsection{Gradient}

Define:
\[
G_\psi(q)
=
\nabla_q\Psi^\perp(q)
\in\R^{m\times d}.
\]

Component:
\[
(G_\psi)_{\alpha\ell}
=
\frac{\partial\psi_\alpha^\perp}{\partial q_\ell}.
\]

\subsection{Mori matrix}

\[
A_{\alpha\gamma}
=
\int
\nabla\psi_\alpha^\perp(q)^\top
D_x(q)
\nabla\psi_\gamma^\perp(q)
\,\rho_x(\dd q).
\]

Quadrature:
\[
A
\approx
\sum_i\sum_{\ell}
W_{i\ell}
G_\psi(q_{i\ell})
D_x(q_{i\ell})
G_\psi(q_{i\ell})^\top.
\]

Units:
\[
[A]=\mathrm{mol\,m^{-3}}\cdot \mathrm{m^2\,s^{-1}}\cdot [\psi]^2/[q]^2.
\]

\subsection{Current coupling}

\[
h_{\alpha a}
=
\int
\nabla\psi_\alpha^\perp(q)^\top
D_x(q)
\nabla P_a(q)
\,\rho_x(\dd q).
\]

Quadrature:
\[
h
\approx
\sum_i\sum_\ell
W_{i\ell}
G_\psi(q_{i\ell})
D_x(q_{i\ell})
G_P(q_{i\ell})^\top.
\]

\[
h\in\R^{m\times3}.
\]

\subsection{Mori correction}

Compute pseudoinverse:
\[
A^\#
\]
by eigendecomposition as defined in Section~1.

Then:
\[
C^\psi=h^\top A^\# h.
\]

\[
C^\psi\in\R^{3\times3}.
\]

Symmetrize numerically:
\[
C^\psi\leftarrow\frac12(C^\psi+C^{\psi\top}).
\]

\section{Projected diffusivity-density}
\label{sec:Dproj}

\[
D_{\rm proj}
=
B-C^Q-C^\psi.
\]

Symmetrize:
\[
D_{\rm proj}\leftarrow\frac12(D_{\rm proj}+D_{\rm proj}^{\top}).
\]

If small negative eigenvalues occur from numerical roundoff, project to PSD:

1. Eigendecompose:
\[
D_{\rm proj}=V\diag(\eta_1,\eta_2,\eta_3)V^\top.
\]

2. Set:
\[
\eta_r^+=\max(\eta_r,0).
\]

3. Reconstruct:
\[
D_{\rm proj}^{+}=V\diag(\eta_1^+,\eta_2^+,\eta_3^+)V^\top.
\]

Use this projection only if:
\[
\min_r\eta_r\ge -\tau_D\max_r|\eta_r|,
\]
for example
\[
\tau_D=10^{-10}.
\]

If the negative eigenvalue is larger than tolerance, fail because the primitive
set is physically inconsistent.

\section{Final scalar conductivity}
\label{sec:finalsigma}

The isotropic projected diffusivity-density is:
\[
D_{\rm iso}
=
\frac13\Tr(D_{\rm proj}).
\]

Conductivity:
\[
\sigma_{\rm S/m}
=
\frac{F^2}{RT}D_{\rm iso}.
\]

Here:
\[
F=96485.33212\ \mathrm{C\,mol^{-1}}.
\]

Units:
\[
\frac{F^2}{RT}
\cdot
\mathrm{mol\,m^{-1}\,s^{-1}}
=
\mathrm{S\,m^{-1}}.
\]

Convert:
\[
\sigma_{\rm mS/cm}=10\sigma_{\rm S/m}.
\]

Final result:
\[
\boxed{
\sigma_{\rm mS/cm}
=
10\frac{F^2}{3RT}
\Tr
\left[
\sum_i c_iD_i^{\rm self}
+
\frac12\sum_{i,j}K_{ij}M_{ij}
-
C^Q
-
h^\top A^\#h
\right].
}
\]

\section{Complete numerical algorithm}
\label{sec:algorithm}

Input:
\[
x=(T,\{C_a\},\{\phi_s\},\mathcal L),
\]
where, per Sections~\ref{sec:c2g}--\ref{sec:library}, \(\mathcal L\) supplies
the physical parameter map \(\mathcal P\).

Output:
\[
\sigma_{\rm mS/cm}(x).
\]

Algorithm:

\begin{enumerate}

\item Build molecule counts from \(\mathcal L\) and \((\{C_a\},\{\phi_s\})\)
(Section~\ref{sec:library}), then build coordinate space \(\Omega_x\) and
coordinate vector \(q\) (Section~\ref{sec:configspace}).

\item Build numerical functions:
\[
U_x(q)\ (\text{Section~\ref{sec:forcefield}}),
\qquad
D_x(q)\ (\text{Section~\ref{sec:mobility}}),
\qquad
P_x(q),
\qquad
G_P(q)=\nabla P_x(q)\ (\text{Section~\ref{sec:polarization}}).
\]

\item Build basin set:
\[
A_1,\ldots,A_N
\]
(Section~\ref{sec:basins}).

\item Assign stoichiometry:
\[
\nu_{i,a}.
\]

\item Generate quadrature:
\[
(q_{i\ell},w_{i\ell})
\quad
\text{for each basin }A_i
\]
(Section~\ref{sec:quadrature}).

\item Solve chemical potentials \(\lambda^\star\) by damped Newton
(Section~\ref{sec:massbalance}):
\[
r_a(\lambda)=\sum_i\nu_{i,a}C^\circ Z_i(\lambda)-C_a,
\]
\[
J_{ab}(\lambda)=\sum_i\nu_{i,a}\nu_{i,b}c_i(\lambda),
\]
\[
J\Delta\lambda=-r.
\]
Backtrack until normalized residual decreases. Stop when
\[
\max_a\frac{|r_a|}{\max(C_a,C_{\min})}<10^{-10}.
\]

\item Compute density weights:
\[
W_{i\ell}
=
C^\circ
w_{i\ell}
\exp
\left[
-\frac{U_x(q_{i\ell})}{RT}
+
\sum_a\nu_{i,a}\lambda_a^\star
\right].
\]

\item Compute populations:
\[
c_i=\sum_\ell W_{i\ell}.
\]

\item For every transition edge \((i,j)\), solve the committor PDE
(Section~\ref{sec:fluxes}):
\[
L_xq_{ij}=0,
\quad
q_{ij}|_{A_i}=0,
\quad
q_{ij}|_{A_j}=1.
\]
Use finite elements or finite volumes. The linear system is:
\[
S u=f,
\]
with
\[
S_{mn}
=
\int
\nabla\varphi_m^\top D_x\nabla\varphi_n
\,\rho_x(\dd q).
\]

\item Compute capacity flux:
\[
K_{ij}
=
\int
\nabla q_{ij}^\top D_x\nabla q_{ij}
\,\rho_x(\dd q).
\]
Set:
\[
K_{ji}=K_{ij}.
\]

\item Compute generator:
\[
Q_{ij}=K_{ij}/c_i,
\qquad
Q_{ii}=-\sum_{j\ne i}Q_{ij}.
\]

\item For every transition edge, sample local conditioned transition paths
(Section~\ref{sec:tpmoments}) with Doob drift:
\[
b^{(ij)}(q)
=
\nabla\cdot D_x(q)-\frac{1}{RT}D_x(q)\nabla U_x(q)
+
2D_x(q)\nabla\log q_{ij}(q).
\]
For each path, compute:
\[
\Delta P=P(q_{\rm end})-P(q_{\rm start}).
\]

\item Compute transition moments:
\[
d_{ij}
=
\frac{\sum_\ell w_\ell^{\rm path}\Delta P_\ell}
{\sum_\ell w_\ell^{\rm path}},
\qquad
M_{ij}
=
\frac{\sum_\ell w_\ell^{\rm path}\Delta P_\ell\Delta P_\ell^\top}
{\sum_\ell w_\ell^{\rm path}}.
\]
Set:
\[
d_{ji}=-d_{ij},
\qquad
M_{ji}=M_{ij}.
\]

\item Define unbounded projectors \(\Pi_i(q)\) (Section~\ref{sec:selfcurrent}).

\item Compute self-current tensors:
\[
D_i^{\rm self}
=
\frac{1}{c_i}
\sum_\ell
W_{i\ell}
G_P(q_{i\ell})
\Pi_i(q_{i\ell})
D_x(q_{i\ell})
\Pi_i(q_{i\ell})^\top
G_P(q_{i\ell})^\top.
\]

\item Compute direct tensor:
\[
B=
\sum_i c_iD_i^{\rm self}
+
\frac12\sum_{i,j}K_{ij}M_{ij}.
\]

\item Compute drift:
\[
b_i=\sum_jQ_{ij}d_{ij}.
\]

\item For each Cartesian axis \(a\), solve (Section~\ref{sec:finitestate}):
\[
\begin{bmatrix}
-Q & \one\\
c^\top & 0
\end{bmatrix}
\begin{bmatrix}
\chi_a\\
\lambda_a^Q
\end{bmatrix}
=
\begin{bmatrix}
b_a\\
0
\end{bmatrix}.
\]

\item Compute finite-state correction:
\[
C^Q_{ab}=\sum_i c_i b_{i,a}\chi_{i,b}.
\]

\item Build memory coordinates \(\Psi(q)\) (Section~\ref{sec:mori}).

\item Orthogonalize memory coordinates by state:
\[
\Psi^\perp(q)=\Psi(q)-\bar\Psi_{s(q)}.
\]

\item Compute memory gradient:
\[
G_\psi(q)=\nabla\Psi^\perp(q).
\]

\item Compute Mori matrix:
\[
A=
\sum_i\sum_\ell
W_{i\ell}
G_\psi(q_{i\ell})
D_x(q_{i\ell})
G_\psi(q_{i\ell})^\top.
\]

\item Compute current coupling:
\[
h=
\sum_i\sum_\ell
W_{i\ell}
G_\psi(q_{i\ell})
D_x(q_{i\ell})
G_P(q_{i\ell})^\top.
\]

\item Compute pseudoinverse \(A^\#\) by eigendecomposition.

\item Compute continuous Mori correction:
\[
C^\psi=h^\top A^\#h.
\]

\item Compute projected diffusivity-density:
\[
D_{\rm proj}=B-C^Q-C^\psi.
\]

\item Compute scalar conductivity:
\[
\sigma_{\rm mS/cm}
=
10\frac{F^2}{3RT}\Tr(D_{\rm proj}).
\]

\end{enumerate}

\section{Effect identification}

An effect is any composition, coordinate, or parameter perturbation \(\theta\)
that changes at least one primitive:
\[
c,\ K,\ Q,\ d,\ M,\ D^{\rm self},\ A,\ h.
\]

To identify the effect numerically, perturb \(\theta\) by a small amount
\[
\delta\theta.
\]

Compute central differences:
\[
\partial_\theta X
\approx
\frac{X(\theta+\delta\theta)-X(\theta-\delta\theta)}
{2\delta\theta}
\]
for
\[
X\in
\{c,K,d,M,D^{\rm self},A,h,\sigma\}.
\]

Primitive ownership scores:
\[
S_c(\theta)=\|\partial_\theta c\|_2,
\]
\[
S_K(\theta)=\|\partial_\theta K\|_F,
\]
\[
S_{dM}(\theta)=
\|\partial_\theta d\|_F+\|\partial_\theta M\|_F,
\]
\[
S_{D}(\theta)=\|\partial_\theta D^{\rm self}\|_F,
\]
\[
S_{Ah}(\theta)=\|\partial_\theta A\|_F+\|\partial_\theta h\|_F.
\]

Assign the effect to the largest nonzero ownership score.

\section{Conductivity-effect ledger}

\[
\begin{array}{lll}
\toprule
\text{Effect} & \text{Primitive location} & \text{Computation}\\
\midrule
\text{Free-ion fraction} & c_i & \partial c_i/\partial\theta\\
\text{Ion association} & c_i,K_{ij} & \partial c,\partial K\\
\text{SSIP/CIP balance} & c_i,K_{ij},D_i^{\rm self} & \partial c,\partial K,\partial D\\
\text{Aggregate formation} & c_i,K_{ij},M_{ij},D_i^{\rm self} & \partial c,\partial K,\partial M,\partial D\\
\text{Neutral ligand coordination} & c_i,K_{ij},A,h & \partial c,\partial K,\partial A,\partial h\\
\text{Solvation competition} & U_x\to c_i,K_{ij} & \partial U\to\partial c,\partial K\\
\text{Coulomb association} & U_x\to c_i,K_{ij} & \partial U\to\partial c,\partial K\\
\text{Born/desolvation} & U_x\to c_i,K_{ij} & \partial U\to\partial c,\partial K\\
\text{Activity/nonideality} & U_x\to c_i,K_{ij} & \partial U\to\partial c,\partial K\\
\text{Viscosity} & D_x\to D_i^{\rm self},K_{ij},A & \partial D\to\partial D^{\rm self},\partial K,\partial A\\
\text{Hydrodynamic size} & D_x\to D_i^{\rm self},K_{ij} & \partial D\\
\text{Anion shape/denticity} & U_x,D_x,K_{ij},A,h & \partial U,\partial D,\partial K,\partial A,\partial h\\
\text{Charge-cloud radius} & U_x,D_x,A,h & \partial U,\partial D,\partial A,\partial h\\
\text{Local crowding/free volume} & U_x,D_x,K_{ij},A,h & \partial U,\partial D,\partial K,\partial A,\partial h\\
\text{Li--anion co-motion} & D_i^{\rm self},M_{ij} & D_{\rm Li,A}>0\\
\text{Li--anion anti-correlation} & D_i^{\rm self},M_{ij} & D_{\rm Li,A}<0\\
\text{Partner switching} & K_{ij},d_{ij},M_{ij},A,h & \partial K,\partial d,\partial M,\partial A,\partial h\\
\text{Identity diffusion} & K_{ij},d_{ij},M_{ij} & \partial K,\partial d,\partial M\\
\text{Structural hopping} & K_{ij},d_{ij},M_{ij} & \partial K,\partial d,\partial M\\
\text{Bound neutral polarization} & A,h & \partial A,\partial h\\
\text{Cage/backjump} & C^Q \text{ or } A,h & \partial C^Q,\partial A,\partial h\\
\text{Ion-atmosphere relaxation} & A,h & \partial A,\partial h\\
\text{Electrophoretic drag} & D_x \text{ or } A,h & \partial D,\partial A,\partial h\\
\text{Debye screening} & U_x,D_x,A,h & \partial U,\partial D,\partial A,\partial h\\
\text{Solvent-shell residence} & c_i,K_{ij},A,h & \partial c,\partial K,\partial A,\partial h\\
\text{Ligand-shell residence} & c_i,K_{ij},A,h & \partial c,\partial K,\partial A,\partial h\\
\text{Temperature} & RT,U_x,D_x,K_{ij},F^2/(RT) & \partial_T\sigma\\
\text{Dielectric response} & U_x,D_x,A,h & \partial U,\partial D,\partial A,\partial h\\
\text{Salt concentration} & c_i,U_x,D_x,K_{ij},A,h & \partial c,\partial U,\partial D,\partial K,\partial A,\partial h\\
\text{Additive microviscosity} & D_x,K_{ij} & \partial D,\partial K\\
\text{Bulky-anion/ligand conductivity increase} & D_i^{\rm self}\text{ or }M_{ij} &
D_{\rm Li,A|i}<0\text{ or partner-switch }M_{ij}\\
\bottomrule
\end{array}
\]

\section{Basis refinement without production Green--Kubo}

Let the current memory basis be
\[
\Phi=(\phi_1,\ldots,\phi_m).
\]

For a candidate missing coordinate \(g(q)\), compute:
\[
A_{gg}
=
\int
\nabla g^\top D_x\nabla g\,\rho_x(\dd q),
\]
\[
A_{g\Phi}
=
\int
\nabla g^\top D_x\nabla\Phi\,\rho_x(\dd q),
\]
\[
h_g
=
\int
\nabla g^\top D_xG_P^\top\,\rho_x(\dd q).
\]

Let
\[
A_{\Phi\Phi}
=
\int
\nabla\Phi^\top D_x\nabla\Phi\,\rho_x(\dd q),
\]
\[
h_\Phi
=
\int
\nabla\Phi^\top D_xG_P^\top\,\rho_x(\dd q).
\]

Remove what is already represented:
\[
\tilde h_g
=
h_g-A_{g\Phi}A_{\Phi\Phi}^\#h_\Phi.
\]
\[
\tilde A_{gg}
=
A_{gg}-A_{g\Phi}A_{\Phi\Phi}^\#A_{\Phi g}.
\]

Score:
\[
S(g)
=
\frac{\|\tilde h_g\|_2^2}{\tilde A_{gg}}.
\]

Add:
\[
g^\star=\argmax_g S(g).
\]

Stop when:
\[
\max_gS(g)<\tau_{\rm basis}
\]
and
\[
|\sigma_{n+1}-\sigma_n|<\tau_\sigma.
\]

This is the production convergence procedure. It uses the generator integrals,
not a long-time GK integral.

\section{Convergence proof}

Let the exact conductivity be
\[
\sigma_{\rm GK}
=
\frac{F^2}{3RT}
\sum_{a=1}^3
\langle J_a,(-L_x)^\#J_a\rangle_{\mu_x}.
\]

Let \(V_n\) be the span of all state indicators and memory coordinates used at
basis level \(n\) (i.e.\ the span growing by successive application of the
basis-refinement procedure above). Let \(u_a^\star\) solve:
\[
-L_xu_a^\star=J_a.
\]

Let \(u_{n,a}\in V_n\) solve the Galerkin projection:
\[
\langle v,(-L_x)u_{n,a}\rangle_{\mu_x}
=
\langle v,J_a\rangle_{\mu_x}
\qquad
\forall v\in V_n.
\]

Then:
\[
\sigma_{\rm GK}-\sigma_n
=
\frac{F^2}{3RT}
\sum_{a=1}^3
\|u_a^\star-u_{n,a}\|_E^2,
\]
where
\[
\|f\|_E^2=\langle f,(-L_x)f\rangle_{\mu_x}.
\]

If
\[
\overline{\bigcup_n V_n}^{\|\cdot\|_E}
\]
contains every \(u_a^\star\), then:
\[
\sigma_n\to\sigma_{\rm GK}.
\]

This proves completeness: the projected estimator of
Section~\ref{sec:finalsigma}, refined by the basis-selection procedure above,
converges to the exact Green--Kubo conductivity even though no long-time
Green--Kubo integral is ever evaluated in production.

\section{Conductivity-behavior test suite}
\label{sec:tests}

This section defines tests that a correct implementation must pass or
\emph{explicitly} flag. Each test identifies \emph{which} primitive
(\(c,K,Q,d,M,D^{\rm self},A,h\)) must carry a physical effect.
If an implementation cannot attribute an effect to a specific primitive,
the effect is not encoded.

\subsection{A. Current-correlation and readout tests}

\paragraph{A1: Nernst--Einstein carrier.}
Construct a model with one charged diffusive state \(i\) with diffusivity
\(D\) and charge \(z\). Assert:
\[
\sigma
=
\frac{F^2}{RT}c_i z^2 D.
\]
This tests that \(c_i D_i^{\rm self}\) is the direct carrier term.

\paragraph{A2: Poisson/Mori corrector on a two-state jump.}
Construct a correlated two-state jump with states \(i,j\), flux \(K_{ij}\),
and opposite displacements \(d_{ji}=-d_{ij}\neq0\).
Check that:
\[
\sigma_{\rm direct}
=
\frac{F^2}{RT}\tfrac{1}{2}(K_{ij}|d_{ij}|^2+K_{ji}|d_{ji}|^2),
\]
and that the finite-state Poisson corrector \(C^Q\neq0\) when the system is
persistent, i.e.\ when the state-conditioned drift \(b_i\neq0\).
A direct-only readout is \emph{wrong} for correlated jumps.

\paragraph{A3: Opposite-charge velocity covariance.}
In a two-ion system with charges \(+1,-1\) and negative Li--anion velocity
covariance:
\[
D_{\rm Li,A}<0.
\]
Assert:
\[
D_Q = D_{\rm Li}+D_{\rm A}-2D_{\rm Li,A}>D_{\rm Li}+D_{\rm A}.
\]
Conductivity must \emph{increase} relative to uncorrelated motion.

\paragraph{A4: State-current drift coupling.}
Construct a state with nonzero state-conditioned drift \(b_i\neq0\) (by
choosing \(Q_{ij}d_{ij}\neq0\) for some neighbor \(j\)). Assert \(C^Q\neq0\).
If the only route to a corrector is a long-time GK integral, fail.

\paragraph{A5: Null current mode.}
Construct a memory coordinate \(\psi\) with \(A_{\psi\psi}=0\) but
\(h_\psi\neq0\). The model must fail loudly (nullspace consistency check
of Section~1) and not return a pseudo-inverse correction.

\paragraph{A6: Zero-displacement motif exchange.}
Construct two states \(i,j\) that are identical in composition except
by solvent orientation label. Set \(d_{ij}=0\). Assert \(K_{ij}d_{ij}=0\)
contributes nothing to the direct tensor \(B\). Motif exchange alone
does not produce conductivity.

\paragraph{A7: Structural hop with charge displacement.}
Construct a partner-switch event \(i\to j\) where Li trades anion partner.
Assert \(d_{ij}\neq0\) and that this contributes \(K_{ij}M_{ij}\) to \(B\).
Verify the reverse event has \(d_{ji}=-d_{ij}\).

\paragraph{A8: Detailed balance and stationarity.}
For every generator \(Q\), assert:
\[
\sum_j Q_{ij}=0,
\qquad
c_i Q_{ij}=c_j Q_{ji},
\qquad
\sum_i c_i Q_{ij}=0.
\]
Finite \(\sigma\) is invalid if any of these fail.

\subsection{B. Speciation and population tests}

\paragraph{B1: Free-ion fraction.}
Increase the Coulomb association strength (e.g.\ scale the Li--anion pairing
free energy by \(\alpha>1\)). Assert \(c_{\rm free}\) decreases and
\(c_{\rm CIP}+c_{\rm SSIP}\) increases.

\paragraph{B2: Shared Li pool for mixed salts.}
Supply a formulation with LiPF6 and LiFSI. Assert that Li appears in exactly
one conserved component \(C_{\rm Li}\), not one per salt.
Check \(C_{\rm PF6}\) and \(C_{\rm FSI}\) are separate components.
Check that \(C_{\rm Li}=C_{\rm LiPF6}^{\rm SI}+C_{\rm LiFSI}^{\rm SI}\) after conversion.

\paragraph{B3: SSIP population.}
In a model with CIP and SSIP basins, assert that \(c_{\rm SSIP}\) changes
when solvent donor number or Li--anion separation threshold changes.
Assert that changes in \(c_{\rm SSIP}\) propagate to \(D_{\rm SSIP}^{\rm self}\)
and \(K_{{\rm SSIP},j}\).

\paragraph{B4: Li-additive coordination motif.}
Supply FEC or VC as a neutral additive.
Assert that a Li--FEC coordinated motif basin \(A_{\rm LiL}\) is populated
when the coordination affinity \(\Delta g_{\rm Li,FEC}^{\rm coord}\) is attractive.
If \(c_{\rm LiL}=0\) despite attractive coordination, fail.
A test that finds FEC changes only \(\eta\) and never \(c_{\rm LiL}\) must fail.

\paragraph{B5: Aggregate formation.}
For high salt concentration, assert that aggregate basins (Li$_2$A$^+$,
LiA$_2^-$, Li$_2$A$_2$, higher) appear in \(\{c_i\}\) when the free energies
support them.
Generic stoichiometric clusters must be expressible; no species-name shortcuts.

\paragraph{B6: Neutral cluster charge covariance.}
For a neutral Li$_2$A$_2$ cluster, assert that the model uses
\(z_i^\top D_i^{\rm centers}z_i\) with \(z_i=(+1,+1,-1,-1)^\top\),
not \(z_{\rm net}^2 D_{\rm COM}=0\).

\paragraph{B7: Activity and nonideality.}
Increase ionic strength by scaling \(C_{\rm Li}\). Assert that activity
corrections change \(c_i\) (basin populations) via \(U_x\), not via a terminal
multiplication of the final \(\sigma\).

\paragraph{B8: Born desolvation.}
Change the Born radius \(r_{\rm Born}\) of the anion. Assert that pair and
cluster free energies change and therefore \(c_{\rm CIP}/c_{\rm SSIP}\) changes.

\subsection{C. State-local mobility and obstruction tests}

\paragraph{C1: Viscosity.}
Increase mixture viscosity \(\eta\) by a factor \(\alpha\). Assert that
\(D_x(q)\) scales as \(1/\alpha\) (Stokes--Einstein), and therefore
\(D_i^{\rm self}\), \(K_{ij}\), and \(A\) all change upstream of \(\sigma\).

\paragraph{C2: Temperature.}
Change \(T\). Assert that \(k_BT\) in \(D_x(q)\), the Boltzmann factor
\(\exp[-U_x/(RT)]\), the viscosity (if temperature-dependent), and the
prefactor \(F^2/(RT)\) all respond consistently.

\paragraph{C3: Hydrodynamic radius.}
Increase the hydrodynamic radius \(a_a\) of a carrier. Assert that
\(\zeta_a=6\pi\eta a_a\) increases, \(D_x\) decreases, and \(D_i^{\rm self}\)
decreases proportionally.

\paragraph{C4: Cluster hydrodynamic radius.}
For a charged cluster (CIP, aggregate), assert that \(D_i^{\rm self}\) uses the
\emph{cluster} hydrodynamic radius, not the free-ion radius.

\paragraph{C5: Anion shape and denticity.}
Introduce a bulky or asymmetric anion via its hydrodynamic radius, charge-cloud
radius, and ligand-field asymmetry descriptors (not by name).
Assert that \(D_i^{\rm self}\), \(K_{ij}\), \(A\), and \(h\) all change.
A test row that finds the bulky anion changes only a scalar radius and
not shape/orientation/covariance descriptors must fail.

\paragraph{C6: Free volume and crowding.}
Increase the occupied volume fraction \(\varphi\). Assert that
\(f_{\rm free\ volume}(\varphi)\) increases \(\eta\), which decreases
\(D_x\) and propagates to \(D_i^{\rm self}\) and \(K_{ij}\).
Near maximum packing \(\varphi\to\varphi_{\max}\), the model must fail or
strongly suppress mobility, not extrapolate silently.

\paragraph{C7: Dense high-salt obstruction.}
At high salt molarity (e.g.\ 2--4 M), assert that finite ionic volume is
wired into \emph{both} mobility \(D_x\) and the finite generator rates
\(K_{ij}\). A model that increases salt without changing mobility or rates
must fail.

\subsection{D. Solvent/additive-specific but name-blind tests}

\paragraph{D1: FEC as coordinating neutral ligand.}
Supply FEC via its descriptors (donor number, coordination affinity,
H-bond donor count, molecular volume, polarizability, etc.) without
using the label ``FEC'' inside the generator. Assert:
\begin{enumerate}
\item Li--FEC coordination motif population \(c_{\rm LiFEC}>0\) when
coordination affinity is attractive.
\item The motif participates in \(K_{ij}\) and \(D_i^{\rm self}\).
\end{enumerate}

\paragraph{D2: FEC viscosity and solvation effects are separate.}
Assert that FEC loading enters:
\begin{enumerate}
\item \(\eta\) via \(f_{\rm additive}\) (viscosity path).
\item \(c_{\rm LiFEC}\) via coordination free energy (speciation path).
\end{enumerate}
Both must be non-zero and separately attributable via primitive ownership scores.

\paragraph{D3: FEC + bulky anion anticorrelation.}
Construct a model with moderate FEC loading and a bulky/delocalized anion.
Assert that the state-local covariance
\[
D_Q(\alpha)
=
D_{{\rm Li}|\alpha}+D_{{\rm A}|\alpha}-2D_{{\rm Li,A}|\alpha}
\]
in the Li--FEC--anion SSIP basin is \emph{different} from that in the
ordinary SSIP basin. Assert that the sign of \(D_{{\rm Li,A}|\alpha}\) can
change, and that conductivity can \emph{increase} despite the viscosity penalty.
A model that can only express ``FEC raises \(\eta\)'' and ``bulky anion adds drag''
but cannot express the changed Li--anion cross-current covariance must fail.

\paragraph{D4: Additive-separated SSIP basin.}
Assert that a Li--FEC--anion basin (Li and anion separated by FEC in the
first shell) is representable as a distinct basin \(A_{\rm LiL\cdots A}\)
(Section~\ref{sec:basins}) with its own \(c_i,D_i^{\rm self},K_{ij}\).

\paragraph{D5: No species-name generator keys.}
Inspect all basin labels, motif labels, memory coordinate definitions, and
transition edges. Assert that none contains the string identifiers
``FEC'', ``VC'', ``PF6'', ``FSI'', ``TFSI'', ``EC'', ``DMC'' as
generator keys. These are only allowed as registry lookup labels during
the descriptor-building phase.

\subsection{E. Ion-atmosphere and electrostatic memory tests}

\paragraph{E1: Debye screening length.}
Change ionic strength \(I=\tfrac{1}{2}\sum_i z_i^2 c_i\). Assert that
\(\kappa^2 \propto I\) changes, that \(R_{\rm atmosphere}\) changes,
and that this propagates to \(A\) and \(h\).

\paragraph{E2: Electrophoretic and relaxation drag.}
Assert that the atmosphere resistance has two separately attributable
components:
\[
D_{\rm full} = D_{\rm no\text{-}atm} - \Delta D_{\rm ep} - \Delta D_{\rm rel}.
\]
Report diagonal Li-center drag, diagonal anion-center drag, and
Li--anion cross-relaxation drag independently.

\paragraph{E3: Charge-cloud form factor.}
Increase the charge-cloud radius of the anion. Assert that the atmosphere
resistance kernel \(\Gamma(k;\kappa,\eta,\epsilon_r)\) changes via the
form factor attenuation, reducing relaxation drag.

\paragraph{E4: Neutral-pair atmosphere cancellation.}
Construct a co-located neutral Li--anion pair. Assert that long-wavelength
atmosphere resistance nearly cancels (opposing charge currents cancel in
\(\rho_z(k\to0)\)).

\paragraph{E5: Atmosphere memory capture.}
Slower state exit time (smaller \(K_{ij}\)) relative to atmosphere relaxation
time must increase the atmosphere memory in \(A,h\) and reduce effective
conductivity. This must not be a terminal scalar gate but appear as a
change in the \(A,h\) primitives.

\subsection{F. Multi-center transport tests}

\paragraph{F1: Charge covariance, not net-charge COM.}
For every multi-site state, assert that
\[
D_Q(i) = z_i^\top D_i^{\rm centers} z_i
\]
is used, not \(z_{\rm net}^2 D_{\rm COM}\).
When \(z_{\rm net}=0\), the former can be non-zero; the latter is always zero.

\paragraph{F2: Perfect comotion limit.}
In a Li--anion pair state with equal and perfectly correlated diffusivities
\(D_{\rm Li}=D_{\rm A}=D_0\) and \(D_{\rm Li,A}=D_0\), assert:
\[
D_Q = D_0+D_0-2D_0=0.
\]
This must be the non-degrading completion test for charge-neutral co-motion.

\paragraph{F3: SSIP relative motion.}
In the SSIP basin, assert that the Li--anion cross-mobility
\(D_{\rm Li,A|\text{SSIP}}\) is distinct from its value in the CIP basin.
The relative motion is a state-derived conditional covariance, not a
formal-charge constant.

\paragraph{F4: Bridge/network motif.}
Supply a multidentate anion with coordination sites \(n_s>1\). Assert that
bridge/network basin concentrations are non-zero when anion concentration
and coordination geometry support them. A model that omits bridge/network
populations when multidentate descriptors are present must fail.

\subsection{G. Cage, backjump, and structural memory tests}

\paragraph{G1: Solvent cage trapping.}
Introduce a caged substate with restricted mobility \(\Pi_{\rm caged}=0\).
Assert that \(c_{\rm caged}D_{\rm caged}^{\rm self}=0\) and therefore
cage-trapped population reduces effective conductivity through speciation.

\paragraph{G2: Backjump anticorrelation.}
Construct a two-edge cycle \(i\to j\to i\) with anti-persistent charge
displacements \(d_{ij}=d_{ji}\) (same direction) but opposite physical
expectation (backjump). Assert that \(C^Q>0\) and reduces \(\sigma\)
below the uncorrected \(B\).

\paragraph{G3: Orientation relaxation memory.}
Introduce anion orientation memory coordinate \(\psi_{\rm orient}\).
Assert that changing orientation relaxation rate changes \(A\) and \(h\),
not the direct \(B\), i.e.\ orientation relaxation is a memory correction,
not a direct carrier effect.

\paragraph{G4: Structural cage exchange (zero displacement).}
Construct a mobile/caged exchange event with \(d_{ij}=0\). Assert it
contributes zero to \(K_{ij}d_{ij}\) and appears only through changing
\(c_i\) and therefore state-local \(D_i^{\rm self}\) weights.

\subsection{H. Concentration and high-loading tests}

\paragraph{H1: Salt molarity sweep.}
Sweep salt molarity from 0.1 to 4 M. Assert that:
\begin{enumerate}
\item Carrier density \(c_{\rm free}\) first increases then decreases.
\item Ionic strength changes \(\kappa\) and \(R_{\rm atmosphere}\).
\item Viscosity increases via \(f_{\rm salt}(I)\).
\item Packing \(\varphi\) increases and enters \(D_x\).
\item Activity changes association equilibrium via \(U_x\).
\item All five effects appear separately in the primitive ownership scores.
\end{enumerate}

\paragraph{H2: Flat or overshooting concentration behavior.}
If a test electrolyte (e.g.\ 1 M LiPF6 in EC:DMC) shows a flat or
overshooting conductivity vs.\ concentration curve relative to data, the
attribution test must identify which primitive family owns the residual:
free-carrier mobility \((D_i^{\rm self})\), association \((K_{ij},c_i)\),
atmosphere \((A,h)\), cage/backjump \((C^Q)\), or finite-volume obstruction
\((D_x)\).

\paragraph{H3: Packing saturation.}
Increase volume fraction \(\varphi\) toward \(\varphi_{\max}\). Assert
that \(U_{\rm pack}\to\infty\) and \(D_x\to0\) before packing limit is
exceeded. The model must fail loudly if \(\varphi>\varphi_{\max}\).

\subsection{I. Solvent matrix tests}

\paragraph{I1: Bulk viscosity.}
Increase \(\eta_{\rm bulk}(x)\). Assert that \(\zeta_a=6\pi\eta a_a\)
increases, \(D_x\) decreases, and \(\sigma\) decreases through mobility,
not a terminal correction.

\paragraph{I2: Bulk dielectric.}
Change \(\epsilon_r\). Assert changes in \(\kappa\),
\(U_{\rm Coul}\) (through \(\epsilon_\infty\)), Born solvation, and
\(R_{\rm atmosphere}\). These must change \(c_i\) and \(K_{ij}\) upstream.

\paragraph{I3: Solvent composition.}
Change EC:DMC ratio. Assert that viscosity, dielectric, donor/acceptor
environment, packing, and coordination all respond via their respective
descriptors. Species names must not appear as generator keys.

\paragraph{I4: Temperature dependence.}
Assert joint consistency: \(F^2/(RT)\), \(k_BT\) in \(D_x=k_BTR_x^\#\),
Boltzmann weights \(\exp[-U_x/(RT)]\), and viscosity \(\eta(T)\) all change
monotonically and consistently with direction of temperature change.

\subsection{J. Audit and dataset effects}

\paragraph{J1: Primitive-gap decomposition.}
For any composition, the implementation must be able to report:
\[
\sigma = \sigma_{\rm vehicular} + \sigma_{\rm jump/memory} - \sigma_{\rm correctors},
\]
where:
\[
\sigma_{\rm vehicular}
=
\frac{F^2}{3RT}
\Tr\bigl[\textstyle\sum_i c_iD_i^{\rm self}\bigr],
\]
\[
\sigma_{\rm jump/memory}
=
\frac{F^2}{3RT}
\Tr\bigl[\tfrac12\textstyle\sum_{i,j}K_{ij}M_{ij}\bigr],
\]
\[
\sigma_{\rm correctors}
=
\frac{F^2}{3RT}
\Tr\bigl[C^Q+h^\top A^\#h\bigr].
\]
If the implementation cannot split these, the primitives are not separately
accessible.

\paragraph{J2: Motif population ledger.}
For every formulation, report:
\[
c_{\rm free},\quad c_{\rm SSIP},\quad c_{\rm CIP},
\quad c_{\rm LiL},\quad c_{\rm agg},\quad c_{\rm bridge}.
\]
These must be non-negative and sum to satisfy mass balance.

\paragraph{J3: Top primitive attribution.}
Perturb each primitive family (\(c\), \(K\), \(d/M\), \(D^{\rm self}\),
\(A,h\)) by \(\pm\delta\) and report the primitives in decreasing order of
\(\partial\sigma/\partial X\). The top attribution should be physically
interpretable.

\paragraph{J4: Ablation audit.}
The following ablations are \emph{diagnostic only} (never production corrections):
\begin{enumerate}
\item Association off: set \(c_{\rm CIP}=c_{\rm SSIP}=0\).
\item Atmosphere off: set \(R_{\rm atmosphere}=0\).
\item Viscosity off: set \(\eta=\eta_{\rm bulk}\) only.
\item Obstruction off: set \(f_{\rm free\ volume}=1\).
\item Cage/memory off: set \(C^Q=0\), \(h=0\).
\end{enumerate}
Report the change in \(\sigma\) for each ablation.

\paragraph{J5: No row dropping.}
Failed or non-converged compositions must be recorded with their reason,
not silently removed from metrics or test sets.

\subsection{K. Descriptor completeness tests}

Before computing conductivity, assert that the species library \(\mathcal L\)
supplies the following for every species:

\[
\begin{array}{ll}
\toprule
\text{Descriptor family} & \text{Required fields} \\
\midrule
\text{Identity} & z_{su},\ M_s,\ \rho_s \\
\text{Size/shape} & a_{su}\ (\text{hydrodynamic}),\ \sigma_{su}\ (\text{LJ}),\
v_s\ (\text{molecular volume}) \\
\text{Electrostatics} & \epsilon_{su},\ \alpha_{su},\ r_{\rm Born,s},\
r_{\rm cloud,s} \\
\text{Coordination/solvation} & \Delta g_{aL}^{\rm coord},\ n_s^{\rm sites},\
n_s^{\rm HBA},\ n_s^{\rm HBD} \\
\text{Transport} & \eta_s^{\rm pure},\ \epsilon_{r,s}^{\rm pure},\ T \\
\bottomrule
\end{array}
\]

If any required field is missing and no default is declared, fail before
returning a conductivity value.

\subsection{L. Minimal pass/fail checklist}

An implementation must \textbf{fail} if any of the following are true:

\begin{enumerate}
\item FEC (or any neutral additive with attractive coordination affinity)
changes \(\eta\) but creates no Li-ligand motif concentration
\(c_{\rm LiL}\).
\item A bulky or asymmetric anion changes only a scalar radius and not its
shape, charge-cloud, orientation/denticity, or Li--anion covariance.
\item SSIP relative/anti-correlated motion is hard-coded from formal charge
rather than computed as conditional covariance.
\item A neutral multi-charge state uses \(z_{\rm net}^2 D_{\rm COM}\) instead
of \(z_i^\top D_i^{\rm centers} z_i\).
\item Mixed salts count Li once per salt source, violating the shared-pool
requirement.
\item A pure motif-exchange event \((d_{ij}=0)\) has nonzero contribution to
the direct conductivity tensor \(B\).
\item High salt loading changes carrier count but not finite volume, viscosity,
activity, atmosphere resistance, or mobility simultaneously.
\item Ion atmosphere appears as a scalar terminal correction to \(\sigma\),
not as entries in \(A\) and \(h\).
\item The Poisson/Mori correctors \(C^Q\) and \(h^\top A^\# h\) are both
identically zero in a system that has cage, atmosphere, partner-switch,
or state-current memory.
\item A generated state label or generator key contains any species name
(FEC, VC, PF6, FSI, TFSI, EC, DMC, or similar) after the descriptor-parsing
stage.
\item The audit cannot attribute the conductivity residual to one of
\(c_i,\ Q_{ij},\ d/M,\ D^{\rm self},\ A/h\).
\item The displacement scale check \(\|\Delta P_{ij,\ell}\|\le L_{\max}\)
is not enforced and meter-scale displacements are accepted.
\end{enumerate}

\section{Validation rules for implementation}

A numerical implementation must fail if:
\[
c_i\le0,
\]
\[
K_{ij}<0,
\]
\[
K_{ij}\ne K_{ji},
\]
\[
\sum_jQ_{ij}\ne0,
\]
\[
c_iQ_{ij}\ne c_jQ_{ji},
\]
\[
d_{ji}\ne-d_{ij},
\]
\[
M_{ji}\ne M_{ij},
\]
\[
D_i^{\rm self}\not\succeq0,
\]
\[
A\not\succeq0,
\]
or if
\[
\|\Delta P_{ij,\ell}\|>L_{\max}.
\]

The projected analytical estimator should consume numerical generator inputs
built from the species library \(\mathcal L\) (Sections~\ref{sec:library}--%
\ref{sec:mobility}), or equivalent primitive arrays. A property database row
with only recipe composition \(r\) and empirical conductivity should not be
converted into a fake generator: by the non-identifiability lemma
(Section~\ref{sec:c2g}), \(r\) alone does not determine \(\sigma\). Standalone
estimator code that imports composition/species utilities but still evaluates
against recipe/composition prediction paths, or validators that call
recipe/composition prediction paths directly, are descriptor-surrogate
implementations, not the full projected analytical model defined in this
document.

\section{Final implementation formula}

\[
\boxed{
\sigma_{\rm mS/cm}(x)
=
10\frac{F^2}{3RT}
\Tr
\left[
\sum_i c_iD_i^{\rm self}
+
\frac12\sum_{i,j}K_{ij}M_{ij}
-
C^Q
-
h^\top A^\#h
\right].
}
\]

Every symbol in this formula is computed by the explicit steps above, starting
from the composition \(x=(T,\{C_a\},\{\phi_s\},\mathcal L)\) through the
species library \(\mathcal L\) (Section~\ref{sec:library}), the force field
and mobility tensor (Sections~\ref{sec:forcefield}--\ref{sec:mobility}), the
generator (Section~\ref{sec:generator}), basins and mass balance
(Sections~\ref{sec:basins}--\ref{sec:massbalance}), reactive fluxes and
transition-path moments (Sections~\ref{sec:fluxes}--\ref{sec:tpmoments}),
self-current (Section~\ref{sec:selfcurrent}), and the finite-state and
continuous memory corrections (Sections~\ref{sec:finitestate}--\ref{sec:mori}).

\end{document}
